% ================================================================
% The Evaporating Universe I: A dissipation principle
% Target: Physics of the Dark Universe (Elsevier)
% ================================================================
\documentclass[a4paper,fleqn]{cas-dc}

\usepackage[numbers]{natbib}
\usepackage{amsmath,amssymb}
\usepackage{graphicx}
\usepackage[section]{placeins}
\graphicspath{{../03_FIGURES/}}
\usepackage{bm}
\usepackage{hyperref}

% --- Custom macros ---
\newcommand{\eps}{\varepsilon}
\newcommand{\epsbare}{\varepsilon_{\rm bare}}
\newcommand{\epsIR}{\varepsilon_{\rm IR}}
\newcommand{\ztrans}{z_{\rm trans}}
\newcommand{\fsurv}{f_{\rm surv}}
\newcommand{\Ikin}{\mathcal{I}}
\newcommand{\rhoc}{\rho_{\rm c}}
\newcommand{\rhoL}{\rho_\Lambda}
\newcommand{\rhob}{\rho_{\rm b}}
\newcommand{\Ocdm}{\Omega_{\rm c}}
\newcommand{\OL}{\Omega_\Lambda}
\newcommand{\Om}{\Omega_{\rm m}}
\newcommand{\wobs}{w_{\rm obs}}
\newcommand{\wlam}{w_\Lambda}
\newcommand{\MPl}{M_{\rm Pl}}

\begin{document}
\let\WriteBookmarks\relax
% --- Permissive float placement (figures inline, not at end) ---
\renewcommand{\topfraction}{0.9}
\renewcommand{\bottomfraction}{0.8}
\renewcommand{\textfraction}{0.1}
\renewcommand{\floatpagefraction}{0.7}
\setcounter{topnumber}{4}
\setcounter{bottomnumber}{4}
\setcounter{totalnumber}{8}

% ================================================================
% FRONT MATTER
% ================================================================

\shorttitle{The Evaporating Universe~I}
\shortauthors{R.\ Alvim}

\title [mode = title]{The Evaporating Universe~I: A dissipation
       principle --- Resolving cosmic tensions without additional
       free parameters}

\author[1]{Ricardo Alvim}[orcid=0009-0004-7744-175X]
\cormark[1]
\ead{ricardoalv1m@pm.me}
\credit{Conceptualization, Methodology, Software, Validation,
        Formal analysis, Investigation, Data curation,
        Writing -- original draft, Writing -- review \& editing,
        Visualization}

\affiliation[1]{organization={Independent Researcher},
            city={Rio de Janeiro},
            country={Brazil}}

\cortext[1]{Corresponding author}

\begin{abstract}
The $\Lambda$CDM model faces two persistent tensions:
the Hubble constant measured locally
($H_0 = 73.17 \pm 0.86$~km\,s$^{-1}$\,Mpc$^{-1}$,
SH0ES) disagrees at $5.8\sigma$ with the CMB inference
($67.36 \pm 0.54$), and weak-lensing surveys measure
$S_8$ values ${\sim}\,2.6\sigma$ below the Planck
prediction. We present the Evaporating Universe (EU),
an interacting dark energy model in which a
Gribov--Zwanziger ghost condensate mediates a continuous
CDM$\,{\to}\,$vacuum energy transfer:
$Q = -\lambda\,\varepsilon(z)\,H\,\rhoc$. The
dissipation fraction ($\lambda = 2/3$), logistic exponent
($b = 19/36$), transition redshift
($z_{\rm trans} \approx 5.986$), bare amplitude
($\varepsilon_{\rm bare} \approx 0.0543$), and effective
coupling ($\varepsilon_{\rm IR} \approx 0.0426$) are derived
from Effective Field Theory in curved spacetime; no parameter
is fitted to data ($\Delta k = 0$). The model is
implemented in a modified CLASS Boltzmann solver and
validated through six Cobaya MCMC runs
(${\sim}\,380{,}000$ total samples; all six converge at
$R-1 < 0.01$)
against Planck CMB, DESI DR2 BAO, Pantheon+ SNe,
SH0ES, and DES-Y3 weak lensing. The CDM drain yields
$H_0^{\rm LKI} = 72.71 \pm 0.30$~km\,s$^{-1}$\,Mpc$^{-1}$,
within $0.52\sigma$ of SH0ES. The $S_8$ tension is
traced to a systematic bias in the $\Lambda$CDM lensing
kernel; forward-modeling the EU kernel reduces the
DES-Y3 discrepancy from ${\sim}\,2.6\sigma$ to $1.03\sigma$.
Including DES-Y3 data directly in the MCMC confirms
compatibility at $0.58\sigma$. A $1024^3$-particle
Gadget-4 $N$-body simulation validates the drain
signature in $P(k)$, and model selection yields
$\Delta\text{AIC} = -6.88$. We present ten
zero-parameter predictions for Euclid DR1, LSST Y1, and
DESI DR3, one of which --- the CPL equation-of-state
mirage ($w_0 = -0.855$ with CMB prior) --- already matches
the DESI DR2 measurement ($w_0 = -0.838 \pm 0.055$) at
$0.3\sigma$.
The EU resolves both cosmological tensions
simultaneously, with zero additional free parameters.
\end{abstract}

\begin{keywords}
dark energy \sep interacting dark energy \sep Hubble tension
\sep $S_8$ tension \sep Gribov--Zwanziger \sep ghost condensate
\sep DESI
\end{keywords}

\maketitle

% ================================================================
% I. INTRODUCTION
% ================================================================
\section{Introduction}
\label{sec:intro}

Two persistent discrepancies challenge the standard
$\Lambda$CDM cosmology. The Hubble constant inferred from
the cosmic microwave background (CMB) under
$\Lambda$CDM~\cite{planck2018},
$H_0 = 67.36 \pm 0.54$~km\,s$^{-1}$\,Mpc$^{-1}$,
disagrees at $5.8\sigma$ with the local distance-ladder
measurement of Breuval, Riess et
al.~\cite{breuval2024},
$H_0 = 73.17 \pm 0.86$~km\,s$^{-1}$\,Mpc$^{-1}$.
Independently, the amplitude of matter fluctuations
$S_8 \equiv \sigma_8\sqrt{\Omega_m/0.3}$ measured by
weak-lensing surveys~\cite{des2022} falls below the
Planck $\Lambda$CDM prediction at $\sim\!2.6\sigma$.
These tensions persist across instruments, calibration
methods, and analysis pipelines, suggesting a systematic
origin rather than a statistical fluctuation.

Numerous extensions of $\Lambda$CDM have been proposed,
including early dark energy~\cite{poulin2019_ede},
interacting dark energy~\cite{divalentino2020},
sign-switching models~\cite{akarsu2023}, and the CPL
parametrization~\cite{chevallier2001,linder2003}
recently favored by DESI DR2~\cite{desi_dr2_2025}.
These approaches typically improve one tension at the
cost of the other, or introduce free parameters whose
values must be fitted to data. No published model
resolves both tensions simultaneously with zero
additional free parameters.

This paper presents the Evaporating Universe (EU), an IDE
model in which a Gribov--Zwanziger (GZ) ghost
condensate~\cite{gribov1978,vandersickel2012}
mediates a continuous transfer of cold dark matter
(CDM) energy into the vacuum:
$Q = -\lambda\,\varepsilon(z)\,H\,\rhoc$, where
$\lambda = 2/3$ is the transversality fraction on the
spatial hypersurface and $\varepsilon(z)$ is a logistic
activation function whose parameters are determined by
the heat-kernel threshold, the CKN holographic
bound~\cite{ckn1999}, and the TT projection of the
ghost--graviton vertex. All five model parameters are
derived from Effective Field Theory in curved spacetime;
none are fitted to observational data ($\Delta k = 0$).

The CDM drain modifies the Friedmann equation at the
percent level, producing two observable effects.
The \textit{global kinematic imprint} (GKI) raises the
background expansion rate to
$H_0^{\rm GKI} = 68.89$~km\,s$^{-1}$\,Mpc$^{-1}$.
The \textit{local kinematic imprint} (LKI), arising from
the KBC void outflow~\cite{keenan2013,wu_huterer2017},
boosts the locally measured value to
$H_0^{\rm LKI} = 72.71$~km\,s$^{-1}$\,Mpc$^{-1}$
--- within $0.52\sigma$ of SH0ES.
The growth tension is resolved by identifying a
systematic bias in the weak-lensing pipeline: the
$\Lambda$CDM lensing kernel applied to EU-generated
data underestimates $S_8$ by construction.
Forward-modeling the EU kernel reduces the DES-Y3
tension from ${\sim}\,2.6\sigma$ to $1.03\sigma$; including
DES-Y3 data directly in the MCMC confirms compatibility
at $0.58\sigma$.
Model selection --- comparing both models on identical
datasets with the KBC void correction applied to
both --- yields $\Delta\text{AIC} = -6.88$ favoring the EU.

The paper is organized as follows.
Section~\ref{sec:theory} derives the five EU parameters
from the GZ framework.
Section~\ref{sec:gki_lki} presents the
analytic predictions for $H_0$.
Section~\ref{sec:class_implementation} describes the
CLASS-EU implementation and stability properties.
Section~\ref{sec:mcmc} reports the
MCMC analysis across six convergence runs.
Section~\ref{sec:drain} characterizes the drain
signature in cosmological observables.
Section~\ref{sec:hubble} resolves the Hubble tension
via the GKI/LKI decomposition.
Section~\ref{sec:nbody} validates the drain signature
with $N$-body simulations.
Section~\ref{sec:growth} resolves the growth
tension through lensing kernel forensics.
Section~\ref{sec:predictions} presents ten falsifiable
predictions for Euclid, LSST, and DESI.
Section~\ref{sec:discussion} places the EU in the
landscape of beyond-$\Lambda$CDM models, discusses
theoretical limitations, and outlines observational
tests. Section~\ref{sec:conclusions} summarizes
the results.

% ================================================================
% II. THEORETICAL FRAMEWORK
% ================================================================
\section{Theoretical Framework}
\label{sec:theory}

We derive the five parameters of the Evaporating Universe from
Effective Field Theory in curved spacetime, utilizing the topological
restrictions of the Gribov--Zwanziger framework and holographic bounds.
No parameter is fitted to observational data ($\Delta k = 0$).

% ----------------------------------------------------------------
% II.A --- The Gribov-Zwanziger Ghost Condensate
% ----------------------------------------------------------------
\subsection{The Gribov--Zwanziger ghost condensate}
\label{sec:gz}

The Gribov--Zwanziger (GZ) framework addresses a fundamental obstacle
in the quantization of non-Abelian gauge theories: the existence of
gauge-equivalent configurations (Gribov copies) that survive the
Landau gauge condition $\partial_\mu A^a_\mu = 0$~\cite{gribov1978}.
The resolution restricts the functional integral to the first Gribov
region,
\begin{equation}
\Omega = \bigl\{A^a_\mu \;\big|\;
  \partial_\mu A^a_\mu = 0,\;\;
  \mathcal{M}^{ab} > 0\bigr\},
\label{eq:gribov_region}
\end{equation}
where $\mathcal{M}^{ab} = -\partial_\mu D^{ab}_\mu$ is the
Faddeev--Popov operator~\cite{vandersickel2012}. This restriction generates a
non-perturbative ghost condensate
$\langle \bar{c}^a c^a \rangle \neq 0$ --- a vacuum expectation
value absent in perturbation theory.

Through Ashtekar's reformulation of general relativity as an SU(2)
gauge theory on the spatial slice
$\Sigma$~\cite{ashtekar1986}, the GZ framework extends to the
gravitational sector~\cite{capri2016}. The ghost condensate on the
Friedmann--Lema\^itre--Robertson--Walker (FLRW) background acts as
the order parameter for a cosmological phase transition: below a
critical redshift $\ztrans \approx 6$, the condensate activates a
continuous transfer of energy from cold dark matter (CDM) to the
vacuum.

CDM participates in this transfer because it is a gauge singlet:
carrying no Standard Model charge, it interacts only gravitationally
and couples to the bulk vacuum through the ghost condensate. Baryonic
matter, bound to $\Sigma$ by its
$\mathrm{SU}(3)_{\rm QCD} \times \mathrm{U}(1)_{\rm EM}$ gauge
charges, is protected from evaporation by those charges.

% ----------------------------------------------------------------
% II.B --- The Unique CDM-Vacuum Interaction
% ----------------------------------------------------------------
\subsection{The unique CDM--vacuum interaction}
\label{sec:interaction}

The modified continuity equations for CDM, vacuum energy, and baryons
are
\begin{align}
\dot{\rhoc} + 3H\rhoc &= -Q, \label{eq:cont_cdm} \\
\dot{\rhoL} &= +Q, \label{eq:cont_de} \\
\dot{\rhob} + 3H\rhob &= 0, \label{eq:cont_b}
\end{align}
where $Q \geq 0$ is the interaction rate. Total energy conservation
is guaranteed by construction:
Eqs.~(\ref{eq:cont_cdm})--(\ref{eq:cont_b}) sum to the Bianchi
identity
$\dot{\rho}_{\rm tot} + 3H(\rho_{\rm tot} + p_{\rm tot}) = 0$.

The functional form of $Q$ is the unique leading-order interaction
consistent with FLRW symmetry and decay kinematics.

\textit{(i)~Dimension.}~From Eq.~(\ref{eq:cont_cdm}),
$[Q] = [\rho]\,[T^{-1}]$: an energy density per unit time.

\textit{(ii)~Initial-state population.}~The interaction describes
CDM$\,\to\,$vacuum transfer. By Fermi's golden rule, the rate is
proportional to the initial-state population: $Q \propto \rhoc$.
This excludes $Q \propto \rhoL$ and mixed forms.

\textit{(iii)~Geometric clock.}~The factor $[T^{-1}]$ must be a
homogeneous, isotropic, first-order geometric scalar. In FLRW
spacetime, the unique candidate is $H = \dot{a}/a$. Higher-order
terms ($\dot{H}/H$, $H^2/H_0$) are sub-leading in the derivative
expansion; particle-physics scales break geometric universality;
the cosmic time $1/t$ is not covariant under time reparametrization.

The unique leading-order interaction is therefore
\begin{equation}
\boxed{Q = \lambda\,\eps(z)\,H\,\rhoc\,,}
\label{eq:Q}
\end{equation}
where $\lambda$ is the dissipation fraction and $\eps(z)$ is the
redshift-dependent coupling amplitude, both derived in
Sec.~\ref{sec:params}. The interacting dark energy (IDE)
literature~\cite{valiviita2008, gavela2009, divalentino2020}
postulates this form; Eq.~(\ref{eq:Q}) establishes it as the unique
solution of the derivative expansion.

The coupling amplitude $\eps(z)$ is the order parameter of the ghost
condensate. Below the condensation threshold $\ztrans$, it follows a
generalized logistic profile with a strict Heaviside cutoff
(Fig.~\ref{fig:anatomy}(a)):
\begin{equation}
\eps(z) = \begin{cases}
\dfrac{\epsIR}{\displaystyle
  1 + \left(\frac{1+z}{1+\ztrans}\right)^{\!1/b}}
  & z \leq \ztrans, \\[10pt]
0 & z > \ztrans.
\end{cases}
\label{eq:eps}
\end{equation}
The cutoff encodes the physics of the condensation: the ghost order
parameter is strictly zero above the threshold, so no CDM is drained
at $z > \ztrans$. The CMB, big bang nucleosynthesis, and all
pre-condensation physics are exactly $\Lambda$CDM. At late times
($z \to 0$), $\eps \to \epsIR$ (full coupling); at the transition
($z = \ztrans$), $\eps = \epsIR/2$ (onset).

% ----------------------------------------------------------------
% II.C --- Framework-Derived Parameters
% ----------------------------------------------------------------
\subsection{Framework-derived parameters}
\label{sec:params}

The five quantities $\{\lambda,\, b,\, \ztrans,\, \epsbare,\, \epsIR\}$
entering Eqs.~(\ref{eq:Q})--(\ref{eq:eps}) are each determined by the
mathematical structure of the GZ framework. We derive them in order of
increasing complexity.

% ---- λ = 2/3 ----
\subsubsection{Dissipation fraction: $\lambda = 2/3$}
\label{sec:lambda}

The dissipation fraction quantifies the fraction of ghost-field
channels available for CDM--vacuum energy transfer. Its value follows
from the transversality structure of the ghost field on the spatial
slice $\Sigma$.

Under the ADM $3{+}1$ decomposition, the ghost field
$c^\mu = (c^0, c^i)$ separates into a constrained temporal component
$c^0$ and three spatial components $c^i$ on $\Sigma$. The Helmholtz
decomposition theorem yields
\begin{equation}
c^i = c^i_T + \partial^i \varphi, \qquad
\partial_i c^i_T = 0,
\label{eq:helmholtz}
\end{equation}
where $c^i_T$ carries two transverse degrees of freedom and $\varphi$
is a longitudinal scalar.

Two independent arguments force $\varphi = 0$. First, the Landau
gauge condition $\partial_\mu A^a_\mu = 0$ projects the
Faddeev--Popov operator onto the transverse functional space,
requiring $\partial_i c^i = 0$ and hence
$\nabla^2\varphi = 0$~\cite{gribov1978}. Second, a longitudinal
condensate
$\langle \bar{c}^a_\parallel c^b_\parallel \rangle \neq 0$ would
single out a preferred direction, spontaneously breaking the isotropy
of the FLRW background. Both routes eliminate the longitudinal mode:
only the two transverse polarizations propagate within $\Omega$.

The GZ effective action confirms this counting. At one-loop order in
$d$ dimensions~\cite{vandersickel2012},
\begin{align}
\Gamma^{(1)}_\gamma &\supset
  \underbrace{-d\,(N^2{-}1)\,\gamma^4}_{\text{tree: }
    d\text{ modes}}
  \notag \\
  &+\;\underbrace{\tfrac{N^2{-}1}{2}(d{-}1)\!\int\!
    \frac{d^dq}{(2\pi)^d}
    \ln\frac{q^4 + 2g^2 N\gamma^4}{q^2}}_{\text{loop: }
    (d{-}1)\text{ transverse modes}}\;,
\label{eq:gz_action}
\end{align}
where the tree-level prefactor $d$ sums over all Lorentz indices and
the one-loop prefactor $(d{-}1)$ reflects the Landau gauge projection.
The gap equation $\partial\Gamma/\partial\gamma^2 = 0$ produces the
ratio $(d{-}1)/d$ as an algebraic quotient of two fixed coefficients
in the action. On the spatial section of FLRW ($d \to n = 3$):
\begin{equation}
\boxed{\lambda = \frac{n-1}{n} = \frac{2}{3}\,.}
\label{eq:lambda}
\end{equation}

The stability of this result under radiative corrections
follows from the Slavnov--Taylor
identity~\cite{taylor1971} for the ghost--gluon vertex in
the Landau gauge:
\begin{equation}
Z_c = Z_g^{-1}\,Z_A^{-1/2}\,,
\label{eq:ward}
\end{equation}
where $Z_c$, $Z_g$, and $Z_A$ are the renormalization
constants of the ghost field, coupling, and gluon field,
respectively. This identity constrains the ghost
self-energy to all orders: the transverse fraction
$(n{-}1)/n$ entering the gap equation is not an
artifact of the one-loop truncation but a consequence
of the gauge structure itself. Dudal
et~al.~\cite{dudal2008} verified that the refined
Gribov--Zwanziger action preserves this identity, so
the channel counting and the gap equation are anchored
to the same algebraic structure. Any higher-loop
correction that modified $\lambda$ would violate
Eq.~(\ref{eq:ward}).

% ---- b = 19/36 ----
\subsubsection{Logistic exponent: $b = 19/36$}
\label{sec:b}

The exponent $b$ in Eq.~(\ref{eq:eps}) controls the sharpness of the
condensation transition. It equals the fraction of the ghost-graviton
interaction that couples to the physical (transverse-traceless)
graviton polarizations in $d = 4$ spacetime.

The ghost-graviton vertex
$V^\mu_{\;\nu;\rho\sigma}(p,q)$ from linearized quantum
gravity~\cite{codello2009} couples the Faddeev--Popov ghost to the
symmetric graviton tensor. The physical content is extracted by
projecting onto the transverse-traceless (TT) sector via the York
projector
$\Pi^{\rm TT}_{\rho\sigma,\lambda\tau}(k)$~\cite{york1973}. The TT
fraction is
\begin{equation}
\boxed{b = \frac{{\rm Tr}\!\left(
  V \cdot \Pi^{\rm TT} \cdot V\right)}{{\rm Tr}\!\left(
  V \cdot \mathbb{I} \cdot V\right)} = \frac{19}{36}\,,}
\label{eq:b}
\end{equation}
verified with irreducible exact rational arithmetic across four
independent momentum configurations (Appendix~\ref{app:tt_vertex}). The
numerator 19 is prime, confirming irreducibility --- a genuine
$d = 4$ tensor invariant.

The TT fraction governs the logistic activation through the
renormalization group (RG). Only the TT-projected component of the
vertex drives the condensate running --- analogous to active quark
flavors modifying the QCD $\beta$-function --- so the ghost
self-energy acquires the anomalous dimension $b$, yielding the Wilson
critical exponent $\nu = 1/b = 36/19$. The logistic $\beta$-function,
\begin{equation}
\frac{d\sigma}{d\ln\mu} =
  -\frac{1}{b}\,\sigma\!\left(1 -
  \frac{\sigma}{\sigma_*}\right),
\label{eq:beta}
\end{equation}
is exact at the upper critical dimension $d = 4$: higher-order
corrections are strictly irrelevant by the Ginzburg criterion,
protecting the universality of the logistic form. The RG scale maps
to redshift through the CKN saturation
(Sec.~\ref{sec:ztrans}):
$\Lambda_{\rm UV}^4 = \rho_{\rm crit} \propto H^2$ implies
$\mu \propto H^{1/2}$, a monotonically increasing function of
$(1{+}z)$. The logistic solution depends only on the ratio
$(1{+}z)/(1{+}\ztrans)$, rendering the result robust against the
precise functional form of $\mu(z)$.

% ---- z_trans ----
\subsubsection{Condensation threshold: $\ztrans \approx 5.986$}
\label{sec:ztrans}

The transition redshift marks the onset of the ghost condensation. It
is determined by a closed chain of five links, each independently
established.

\textit{Link~1: One-loop amplitude.}~The universal one-loop prefactor
from the Seeley--DeWitt heat-kernel expansion~\cite{vassilevich2003}
in $d = 4$ dimensions is
\begin{equation}
(4\pi)^{-d/2}\big|_{d=4} = \frac{1}{16\pi^2} \approx 0.00633\,.
\label{eq:hk_prefactor}
\end{equation}
This is the standard phase-space suppression of a single quantum loop,
appearing identically in dimensional regularization, the
Coleman--Weinberg potential, and the ABJ anomaly.

\textit{Link~2: Holographic bound.}~Cohen, Kaplan, and
Nelson~\cite{ckn1999} established that an effective field theory
in a region of size $L$ must satisfy
$L^3\Lambda_{\rm UV}^4 \lesssim L \MPl^2$. Applied to the Hubble
sphere ($L = H^{-1}$) with the exact Schwarzschild condition:
\begin{equation}
\Lambda_{\rm UV}^4 \leq
  \frac{3\MPl^2 H^2}{8\pi} = \rho_{\rm crit}\,.
\label{eq:ckn}
\end{equation}

\textit{Link~3: Saturation.}~The CKN bound
(Eq.~\ref{eq:ckn}) is formally an inequality. In a spatially
flat universe
($\Omega_k = 0.001 \pm 0.002$~\cite{planck2018}), the
late-time de~Sitter attractor drives
$\Omega_\Lambda \to 1$, saturating the Gibbons--Hawking
entropy bound $S = A/(4G)$ as a thermodynamic
identity~\cite{gibbons_hawking1977}. The CKN energy bound is
the effective field theory formulation of this holographic
constraint~\cite{ckn1999}; the inequality becomes an equality:
$\Lambda_{\rm UV}^4 = \rho_{\rm crit}$. Conversely, exact
saturation requires $\Omega_k = 0$: any spatial curvature
introduces a time-dependent correction
$\propto k/(a^2 H^2)$, breaking the constancy of the
bound. The observed flatness and CKN saturation are
mutually required.

\textit{Link~4: Threshold.}~With
$\Lambda_{\rm UV}^4 = \rho_{\rm crit}$ (Link~3), the
cosmological density parameter
$\Omega_{\rm DE} \equiv \rho_\Lambda/\rho_{\rm crit}$
measures the vacuum energy as a fraction of the total UV
phase space. Combining with Link~1, the condensate
nucleates when this fraction reaches the one-loop
amplitude:
\begin{equation}
\Omega_{\rm DE}(\ztrans) = \frac{1}{16\pi^2}\,.
\label{eq:threshold}
\end{equation}

\textit{Link~5: Solution.}~In the matter-dominated $\Lambda$CDM
background (valid at $z \sim 6$ where EU and radiation corrections
are negligible):
\begin{equation}
\boxed{\ztrans = \left[
  \frac{\OL\,(16\pi^2 - 1)}{\Om}\right]^{1/3} - 1 \approx 5.986\,.}
\label{eq:ztrans}
\end{equation}
The inputs are $\OL = 0.6847$ and $\Om = 0.3153$ from
Planck~2018~\cite{planck2018}. Because $\eps(z) = 0$ for
$z > \ztrans$ (Eq.~\ref{eq:eps}), the pre-condensation background
is rigorously $\Lambda$CDM; Eq.~(\ref{eq:ztrans}) therefore evaluates
exactly using the primordial densities. The sensitivity to
$h = H_0/100$ is weak: $\ztrans$ varies by $\sim 6\%$ over
$h \in [0.66, 0.70]$
(Appendix~\ref{app:ztrans_sensitivity}).

% ---- ε_bare, ε_IR ----
\subsubsection{Coupling amplitude: $\epsbare$ and $\epsIR$}
\label{sec:eps}

The coupling amplitude has two components: a bare value set by the
vacuum channel counting and an infrared screening from the Gribov
horizon.

\textit{Bare coupling.}~The Gribov--Zwanziger framework
applies to any gauge theory formulated in terms of a
connection. General relativity on $\Sigma$, as
conventionally written in metric variables, is not a
gauge theory in this sense. Ashtekar's canonical
reformulation~\cite{ashtekar1986} resolves this by
recasting the gravitational phase space in terms of an
SU(2) connection $A_a^i$ and a densitized triad
$\tilde{E}^a_i$ on the spatial slice. Einstein's
constraints become polynomial in these variables, and
the internal gauge group is SU(2) with
$\dim(\mathrm{SU}(2)) = n = 3$. The existence of Gribov
copies is a topological property of all non-Abelian gauge
theories in the Landau gauge, independent of the specific
gauge group~\cite{vandersickel2012}. This SU(2)
identification therefore places the gravitational sector
on the same formal footing as QCD, making the
Faddeev--Popov operator, the Gribov region $\Omega$, and
its boundary $\partial\Omega$ well-defined.

Within this gauge-theoretic formulation, the GZ
localizing field
$\varphi_a^{ij}$~\cite{zwanziger1989} --- an auxiliary
tensor coupling the ghost ($c^i$) and anti-ghost
($\bar{c}^j$) to the gauge index $a$ --- possesses three
independent indices, each running from 1 to $n$. In the
deep UV, where the ghost and anti-ghost are distinct
fermionic fields with unrestricted quantum numbers
(Zwanziger's independent-index
theorem~\cite{zwanziger1989}), the auxiliary field
fluctuates as a fully general rank-3 tensor:
$N_{\rm bare} = n^3 = 27$ bare channels. In the
classical IR, the metric symmetry
$\gamma_{ij} = \gamma_{ji}$ constrains the ghost
indices, reducing the auxiliary field to
$n\,n(n{+}1)/2$ independent components:
$N_{\rm classical} = 18$ for $n = 3$.

The Gribov projection connects these two regimes.
Setting the effective count
$\lambda\, n^3 = n^2(n{-}1)$ equal to the
classical count:
\begin{equation}
n^2(n{-}1) = \frac{n^2(n{+}1)}{2}
\;\;\Longrightarrow\;\;
\boxed{n = 3\,.}
\label{eq:n3}
\end{equation}
The quantum vacuum recovers the classical limit only in
three spatial dimensions: $(2/3)(27) = 18
= 3(3{+}1)/2 \cdot 3$. This consistency
condition is not imposed; it is a consequence of
requiring the GZ ghost sector to reproduce the correct
number of gravitational degrees of freedom.

\textit{Democratic coupling.}~CDM carries no Standard
Model gauge charge and, crucially, no internal SU(2)
quantum numbers under the Ashtekar connection. It is a
gauge singlet. By the Wigner--Eckart theorem, the
matrix element of any gauge-invariant operator between
singlet states is proportional to the identity on the
mediator space: no channel is preferred over any other.
CDM therefore couples with equal amplitude to all 27
bare channels. The transverse one-loop factor
$(4\pi)^{-1}$ --- reflecting the restriction of the
ghost dynamics to the two-dimensional transverse plane
on $\Sigma$ (Sec.~\ref{sec:lambda}), in contrast to
the bulk factor $(4\pi)^{-2}$ used for the condensation
threshold --- distributes uniformly over the bare
modes, yielding the phase-space fraction
$1/(4\pi\,n^3)$. By Fermi's golden rule, the
transition rate scales as $|\varepsilon|^2$; the
coupling amplitude is therefore the square root of
this fraction:
\begin{equation}
\boxed{\epsbare = \frac{1}{\sqrt{4\pi\, n^3}}
  = \frac{1}{\sqrt{108\pi}} \approx 0.05429\,.}
\label{eq:epsbare}
\end{equation}

\textit{IR screening.}~In the deep infrared, the Gribov horizon
modifies the ghost propagator from
$D_{\rm UV}(k) = 1/(k^2 + \gamma^2)$ to
$D_{\rm GZ}(k) = k^2/(k^4 + \gamma^4)$. This deformation of
the internal propagator acts as a vertex renormalization,
screening the bare coupling by a calculable factor. The
ratio of integrated propagator densities on the
two-dimensional transverse momentum
plane~\cite{zwanziger1991} evaluates exactly to
(Appendix~\ref{app:gribov_integral})
\begin{equation}
\frac{I_{\rm IR}}{I_{\rm UV}} = \frac{\pi}{4}\,.
\label{eq:screening}
\end{equation}
The effective coupling is therefore
\begin{equation}
\boxed{\epsIR = \frac{\pi}{4}\,\epsbare \approx 0.04264\,.}
\label{eq:epsIR}
\end{equation}

% ---- Table I ----
\begin{table}[pos=!htbp]
\caption{The five EU parameters. All are derived from the
  mathematical structure of the Gribov--Zwanziger EFT in curved
  spacetime. None are fitted to observational data ($\Delta k = 0$).}
\label{tab:params}
\begin{tabular}{llll}
Parameter & Value & Exact form & Origin \\
\hline
$\lambda$ & $0.6\overline{6}$ & $2/3$ &
  GZ transversality \\
$b$ & $0.5278$ & $19/36$ &
  TT vertex fraction \\
$\ztrans$ & $5.986$ & Eq.~(\ref{eq:ztrans}) &
  Heat kernel + CKN \\
$\epsbare$ & $0.05429$ & $1/\sqrt{108\pi}$ &
  GZ channels \\
$\epsIR$ & $0.04264$ & $(\pi/4)\,\epsbare$ &
  Gribov screening \\
\end{tabular}
\end{table}

% ----------------------------------------------------------------
% II.D --- Modified Friedmann Equation
% ----------------------------------------------------------------
\subsection{Modified Friedmann equation}
\label{sec:friedmann_eu}

The CDM density evolves under the combined effect of Hubble dilution
and vacuum transfer. Integrating Eq.~(\ref{eq:cont_cdm}):
\begin{equation}
\rhoc(a) = \rho_{c,0}\,a^{-3}\,\fsurv(a),
\label{eq:rho_c}
\end{equation}
where the survival fraction is
\begin{equation}
\fsurv(a) = \exp\!\left[
  -\lambda \int_0^{a}
  \eps(a')\,\frac{da'}{a'}\right].
\label{eq:fsurv}
\end{equation}

The Heaviside cutoff in Eq.~(\ref{eq:eps}) renders the total
accumulated drain analytically tractable. The kinematic integral is
\begin{equation}
\Ikin = \epsIR\,b\left[
  \ln\!\left(1 + (1{+}\ztrans)^{1/b}\right) - \ln 2\right]
  \approx 0.0679\,,
\label{eq:Ikin}
\end{equation}
where the $-\ln 2$ arises from the upper limit of integration at
$\ztrans$. The CDM survival fraction at $z = 0$ is
\begin{equation}
\fsurv(0) = e^{-\lambda\Ikin} \approx 0.9558\,,
\label{eq:fsurv0}
\end{equation}
corresponding to a $4.4\%$ transfer of CDM energy to the vacuum since
$z \approx 6$.

The modified Friedmann equation is
\begin{align}
H^2(a) = H_0^2\big[
  &\Omega_{\rm r}\,a^{-4}
  + \Omega_{\rm b}\,a^{-3}
  \notag \\
  &+ \Ocdm\,a^{-3}\,\fsurv(a)
  + \OL(a)\big],
\label{eq:friedmann}
\end{align}
where $\Omega_{\rm r}$ encodes photons and relativistic neutrinos
(one massive species with $\sum m_\nu = 0.06$~eV~\cite{planck2018}),
$\Ocdm$ denotes the pre-drain CDM density parameter anchored by the
CMB ($\eps = 0$ at $z > \ztrans$, so $\Ocdm$ retains its
$\Lambda$CDM value), and $\OL(a)$ is obtained by integrating
Eq.~(\ref{eq:cont_de}) numerically. The CDM and vacuum terms do not cancel: drained CDM
energy acquires the vacuum equation of state ($w = -1$) and ceases to
dilute, generating a net modification to $H(z)$. The full numerical
integration is performed by the CLASS Boltzmann
solver~\cite{blas2011} (Sec.~\ref{sec:class_implementation}).

Despite absorbing CDM energy, the vacuum retains $w_\Lambda = -1$
exactly --- a property proven algebraically in
Sec.~\ref{sec:stability}.

With all five parameters determined, the model predicts two distinct
kinematic imprints on the Hubble rate: a global shift from CDM drain
and a local boost from the KBC void outflow.

% ================================================================
% III. KINEMATIC IMPRINTS
% ================================================================
\section{Kinematic Imprints: GKI and LKI}
\label{sec:gki_lki}

The five framework-derived parameters produce two distinct kinematic
signatures on the Hubble rate. The \textit{Global Kinematic Imprint}
(GKI) modifies the background expansion through CDM drain, while the
\textit{Local Kinematic Imprint} (LKI) biases local distance-ladder
measurements through the KBC void outflow. Both are zero-parameter
predictions.

% ----------------------------------------------------------------
% III.A --- GKI
% ----------------------------------------------------------------
\subsection{Global Kinematic Imprint}
\label{sec:gki}

The CMB power spectrum fixes the physical CDM density
$\omega_{\rm cdm} \equiv \Ocdm h^2 = 0.1200 \pm 0.0012$~\cite{planck2018}
at $z \sim 1100$, where the EU coupling is strictly zero
($\eps = 0$ for $z > \ztrans$). The CMB therefore measures the same
$\omega_{\rm cdm}$ as $\Lambda$CDM.

If CDM evaporates at late times, $\Ocdm$ decreases while
$\omega_{\rm cdm}$ remains CMB-anchored. The reduced Hubble parameter
$h$ must increase to maintain the constraint:
\begin{equation}
h^{\rm EU} = \frac{h^{\rm Planck}}{\sqrt{\fsurv(0)}}\,,
\label{eq:h_gki}
\end{equation}
yielding the global prediction
\begin{equation}
\boxed{H_0^{\rm GKI}
  = \frac{H_0^{\rm Planck}}{\sqrt{\fsurv(0)}}
  = \frac{67.36}{\sqrt{0.9558}}
  \approx 68.90 \;\text{km}\,\text{s}^{-1}\,\text{Mpc}^{-1}\,.}
\label{eq:H0_gki}
\end{equation}

Because CDM has evaporated, the present-day physical matter density
is $\omega_m(z{=}0) = \omega_b + \omega_{\rm cdm}\,\fsurv(0)
\approx 0.137$. The corresponding matter density parameter is
$\Omega_m^{\rm EU} = \omega_m(z{=}0) / (h^{\rm EU})^2 \approx 0.289$,
compared with $\Omega_m^{\Lambda\rm CDM} = 0.3153$.

This mechanism --- CDM drain reducing $\Ocdm$ and raising $h$ ---
is the same physics exploited by Interacting Dark Energy
models~\cite{divalentino2020}, but with a coupling
fixed by the framework ($\lambda = 2/3$) rather than
fitted to data.

% ----------------------------------------------------------------
% III.B --- LKI
% ----------------------------------------------------------------
\subsection{Local Kinematic Imprint}
\label{sec:lki}

The GKI predicts the \textit{global} expansion rate. Local
distance-ladder measurements are systematically higher because the
observer sits inside the KBC void~\cite{keenan2013}, a local
underdensity extending to $R \sim 300\;\text{Mpc}$ with observed
contrast $\delta_{\rm obs} = -0.46 \pm 0.06$.

The void outflow biases the locally inferred $H_0$
(Wu \& Huterer 2017~\cite{wu_huterer2017}):
\begin{equation}
\frac{\delta H_0}{H_0}
  = -\frac{1}{3}\,\delta_{\rm true}\,f(\Omega_m)\,\Theta\,,
\label{eq:void_bias}
\end{equation}
where $f(\Omega_m)$ is the linear growth rate and $\Theta$ accounts
for the non-linear correction of
Marra et al.~\cite{marra2013}.

Two refinements enter the calculation. First, the observed contrast
$\delta_{\rm obs}$ is inflated by redshift-space distortions
(Haslbauer, Banik \& Kroupa 2020~\cite{haslbauer2020}):
\begin{equation}
\delta_{\rm true}
  = \frac{\delta_{\rm obs}}{1 + f/b_g}
  = \frac{-0.46}{1 + 0.514/1.0}
  \approx -0.304\,,
\label{eq:delta_true}
\end{equation}
where $f = 0.514$ is computed from the exact Riccati
ODE~\cite{linder2005} evaluated in the EU background
(not the $\Omega_m^{0.55}$ approximation)
and $b_g = 1.0$ is the $K$-band galaxy bias. Second, the
non-linear correction $\Theta = 1.064$ amplifies the outflow by
$6.4\%$ relative to linear theory~\cite{marra2013}.

The local prediction is
\begin{align}
H_0^{\rm LKI}
  &= H_0^{\rm GKI} + \delta H_0
  \notag \\
  &= 68.90 + 3.82
  \approx \boxed{72.72\;\text{km}\,\text{s}^{-1}\,\text{Mpc}^{-1}\,.}
\label{eq:H0_lki}
\end{align}

CDM drain also causes adiabatic halo expansion~\cite{gnedin2004},
but its effect on the distance ladder is negligible:
$\Delta R_{\rm halo} \sim 10\;\text{kpc}$ versus
$D_{\rm SN} \sim 40\;\text{Mpc}$ implies
$\Delta H_0 \sim 0.009\;\text{km}\,\text{s}^{-1}\,\text{Mpc}^{-1}$.

\textit{Sensitivity.}~The LKI prediction is robust against the
KBC void uncertainty. Propagating
$\delta_{\rm obs} = -0.46 \pm 0.06$ through the full
RSD-corrected formula yields
$H_0^{\rm LKI} \in [72.19, 73.26]$ at $1\sigma$, with worst-case
tension against SH0ES of $1.14\sigma$.

% [Figure 1 placed in Sec. VII.D where the geometric wall is discussed]

% ----------------------------------------------------------------
% III.C --- All tracers agree
% ----------------------------------------------------------------
\subsection{Observational hierarchy}
\label{sec:hierarchy}

The GKI and LKI channels predict a characteristic
hierarchy: global probes recover $H_0$ consistent with
the Planck value, while local calibrators --- Cepheids,
TRGB, JAGB --- converge near
$H_0^{\rm LKI} \approx 72.7$, as their Rung~3 supernovae
lie within the KBC void.

% [TABLE: Observational hierarchy — 5 tracers]
\begin{table}[pos=!htbp]
\caption{EU predictions for the Hubble constant. The CMB prediction
is the GKI (global); all local tracers receive the same LKI (void)
correction. TRGB/JAGB values from
Freedman et al.~\cite{freedman2024}.
All values in km\,s$^{-1}$\,Mpc$^{-1}$.}
\label{tab:hierarchy}
\begin{tabular}{lcccr}
Tracer & $H_0^{\rm obs}$ & $\sigma$ & $H_0^{\rm EU}$ & Tension \\
\hline
CMB/Planck       & 67.36 & 0.54 & 67.36 & $0.0\sigma$ \\
JAGB/JWST        & 67.80 & 2.72 & 72.72 & $1.8\sigma$ \\
TRGB/JWST        & 68.81 & 2.22 & 72.72 & $1.8\sigma$ \\
TRGB (mixed)     & 70.39 & 1.94 & 72.72 & $1.2\sigma$ \\
SH0ES 2024       & 73.17 & 0.86 & 72.72 & $0.5\sigma$ \\
\end{tabular}
\end{table}

Table~\ref{tab:hierarchy} summarizes the comparison. All five
tracers are consistent with the EU prediction at $< 2\sigma$.
The analytical predictions of this section
are validated by the full Boltzmann integration in
Sec.~\ref{sec:class_implementation}.

The same CDM drain that modifies $H_0$ also suppresses structure
growth by reducing $\Omega_m$ and altering $H(z)$. The
quantitative resolution of the $S_8$ tension is presented in
Sec.~\ref{sec:growth}.

\FloatBarrier
% ================================================================
% IV. NUMERICAL IMPLEMENTATION
% ================================================================
\section{Numerical Implementation}
\label{sec:class_implementation}

We implement the EU cosmology in the Cosmic Linear Anisotropy
Solving System (\textsc{class}
v3.3.4~\cite{blas2011}), verifying that the analytic
predictions of Sec.~\ref{sec:gki_lki} are reproduced by the
full Boltzmann integration. The implementation modifies four
source files; the complete patched source, analysis notebooks,
and MCMC chains are publicly
available~\cite{eu_github}.

% ----------------------------------------------------------------
% IV.A — CLASS-EU solver
% ----------------------------------------------------------------
\subsection{Solver architecture}
\label{sec:solver}

Four patches are applied to the \textsc{class} source. (i)~The
header \texttt{background.h} receives the five EU parameters
($\varepsilon_{\rm IR}$, $\ztrans$, $b$, $\lambda$,
and the perturbation switch). (ii)~The parser \texttt{input.c}
reads these parameters from the configuration file, defaulting
to $\varepsilon_{\rm IR} = 0$ (standard $\Lambda$CDM).
(iii)~The background integrator \texttt{background.c} implements
the coupled ODE system
\begin{align}
\frac{d f_{\rm cdm}}{d\ln a}
  &= -\lambda\,\eps(z)\,f_{\rm cdm}\,,
\label{eq:fcdm_ode}\\[4pt]
\frac{d f_{\rm vac}}{d\ln a}
  &= +\lambda\,\eps(z)\,f_{\rm cdm}\,a^{-3}\,,
\label{eq:fvac_ode}
\end{align}
where both $f_{\rm cdm}$ and $f_{\rm vac}$ are normalized to
$\rho_{c,0} \equiv \Omega_{c,0}\,H_0^2/(8\pi G/3)$.
The system is solved by a fourth-order Runge--Kutta integrator
over 2000 steps in $\ln a \in [-7, 0]$, with initial conditions
$f_{\rm cdm} = 1$, $f_{\rm vac} = 0$.
The CDM density is multiplied by $f_{\rm cdm}(a)$ at each
background evaluation, and the vacuum receives the Friedmann
closure injection $\rho_{c,0}\,f_{\rm vac}(a)$.
(iv)~The perturbation integrator \texttt{perturbations.c}
receives the interacting dark energy (IDE) source terms
of Valiviita, Majerotto \& Maartens (2008)~\cite{valiviita2008}.

The EU and $\Lambda$CDM runs share identical cosmological
parameters (Planck 2018
baseline~\cite{planck2018}). The $\Lambda$CDM run
fixes $H_0 = 67.36\;\text{km}\,\text{s}^{-1}\,\text{Mpc}^{-1}$; the EU run
inputs the same $100\theta_s$ and lets \textsc{class} shoot
for~$H_0$, which is the standard procedure when comparing
models with different late-time physics.

% ----------------------------------------------------------------
% IV.B — Perturbation stability
% ----------------------------------------------------------------
\subsection{Perturbation stability}
\label{sec:stability}

The EU coupling $Q = -\lambda\,\eps(z)\,H\,\rhoc$ is
proportional to the local CDM density. This proportionality
guarantees algebraic cancellation of the IDE perturbation
source terms for $w_\Lambda = -1$:

\textit{Continuity.}~The source term in the CDM overdensity
equation is proportional to $[\delta Q/Q - \delta_c]$.
Since $Q \propto \rhoc$, the fractional perturbation
$\delta Q/Q = \delta_c$, yielding
$[\delta_c - \delta_c] = 0$.

\textit{Euler.}~The momentum drag is proportional to
$[\theta_{\rm de} - \theta_{\rm cdm}]$. For $w = -1$ (vacuum),
$\theta_{\rm de} = 0$; in synchronous gauge,
$\theta_{\rm cdm} = 0$ by construction. Hence
$[0 - 0] = 0$.

Dense and underdense regions evaporate at the same fractional
rate. Structure growth is modified only through the background:
the diluted Poisson source (weaker $\rho_{\rm cdm}$) produces
shallower potential wells.

\textit{Null test.}~We verify this cancellation by running
\textsc{class} twice with identical parameters but the IDE
perturbation switch toggled on and off. The difference in
$\sigma_8$ is
\begin{align}
\Delta\sigma_8
  &= |\sigma_8^{\rm IDE\;on} - \sigma_8^{\rm IDE\;off}|
  \notag \\
  &= 0\quad (\text{exact to machine precision})\,.
\label{eq:null_test}
\end{align}

\textit{Doom factor.}~Standard IDE models with $w \neq -1$
suffer from large-scale instabilities when the ratio
$Q/[(1+w)\rho_{\rm de}]$ diverges~\cite{gavela2009,valiviita2008}.
In the EU, $w_\Lambda = -1$ exactly: the vacuum does not
cluster ($\delta_{\rm de} = 0$) and carries no velocity
perturbations ($\theta_{\rm de} = 0$). The dark energy
Euler equation drops out of the perturbation system,
so the singular ratio never arises. Unlike standard IDE
models where $w$ is a free parameter that must be tuned
away from $-1$ to avoid this instability, the EU derives
$w_\Lambda = -1$ algebraically from the Cancellation
Theorem. The doom factor instability is thus structurally
absent --- a consequence of the model's internal
consistency, not a numerical choice.

% ----------------------------------------------------------------
% IV.C — Background results
% ----------------------------------------------------------------
\subsection{Background evolution}
\label{sec:bg_results}

Table~\ref{tab:background} compares the $z = 0$ background
quantities from the fiducial \textsc{class} runs (Planck 2018
input parameters). The EU CDM drain of $4.42\%$ is confirmed,
and energy conservation is exact:
$|\Omega_{\rm tot} - 1| < 10^{-15}$ (Bianchi identity).

% [TABLE: Background ΛCDM vs EU]
\begin{table}[t]
\caption{Background quantities at $z = 0$ from \textsc{class}.
Both runs use Planck 2018 input parameters; the EU run shoots
for $H_0$ via $\theta_s$ matching.}
\label{tab:background}
\begin{tabular}{lccr}
Quantity & $\Lambda$CDM & EU & Shift \\
\hline
$H_0$ (km\,s$^{-1}$\,Mpc$^{-1}$)      & 67.36 & 68.56 & $+1.8\%$  \\
$\Omega_m$             & 0.315 & 0.293 & $-7.1\%$  \\
$\Omega_\Lambda$       & 0.685 & 0.707 & $+3.2\%$  \\
$f_{\rm cdm}(0)$       & 1.000 & 0.956 & $-4.4\%$  \\
$r_d$ (Mpc)            & 147.06 & 147.06 & $0.00\%$ \\
$\sigma_8$             & 0.811 & 0.829 & $+2.3\%$  \\
$S_8$                  & 0.831 & 0.819 & $-1.4\%$  \\
\end{tabular}
\end{table}

The vacuum equation of state from \textsc{class} reads
$w_\Lambda = -1.003$, consistent with $w = -1$ exact to
within the RK4 integration residual ($\sim 0.3\%$).
The sound horizon $r_d$ is identical between the two models
because the EU coupling is zero at $z > \ztrans \approx 6$,
leaving recombination physics unchanged.
The fiducial $H_0 = 68.56$ in Table~\ref{tab:background}
reflects the background response under fixed Planck
primordial densities; re-evaluation at the MCMC posteriors
(Sec.~\ref{sec:drain}) shifts $H_0$ upward into sub-percent
agreement with the analytic prediction of
Sec.~\ref{sec:gki}.

% ----------------------------------------------------------------
% IV.D — CMB consistency
% ----------------------------------------------------------------
\subsection{CMB power spectrum}
\label{sec:cmb_consistency}

The $C_\ell^{TT}$ power spectra of the two models are
virtually indistinguishable at $\ell > 30$, where the acoustic
peaks are determined by pre-recombination physics
($z > \ztrans$). At $\ell \lesssim 20$, the EU spectrum shows
a modest enhancement due to the late integrated Sachs--Wolfe
(ISW) effect: CDM drain reduces the matter density at $z < 2$,
causing gravitational potentials to decay faster and boosting
the low-$\ell$ power. This $\sim 5\%$ ISW enhancement is a
falsifiable prediction of the model.

With the solver validated at both background and perturbation
level, we proceed to the statistical analysis in
Sec.~\ref{sec:mcmc}.

The development of the CLASS-EU solver patches, Cobaya
configuration files, and Gadget-4 simulation scripts was
supported by an AI coding assistant used for code development,
debugging, and optimization. All numerical outputs were
independently verified by the author through convergence
diagnostics ($R-1 < 0.01$), analytical cross-checks
(Table~\ref{tab:concordance}), and null tests
(Table~\ref{tab:nbody_drain}).

% ================================================================
% V. STATISTICAL ANALYSIS
% ================================================================
\section{Statistical Analysis}
\label{sec:mcmc}

We confront the EU model with cosmological data via Markov
Chain Monte Carlo (MCMC) sampling, using
Cobaya~\cite{torrado2021} as the sampler and CLASS-EU
(Sec.~\ref{sec:class_implementation}) as the theory engine.

% ----------------------------------------------------------------
% V.A — Methodology
% ----------------------------------------------------------------
\subsection{Sampling strategy}
\label{sec:methodology}

The analysis comprises six MCMC runs organized in three
phases across two parameter configurations.

The first configuration, \textit{EU-blind} ($\Delta k = 3$),
samples the three EU parameters
($\varepsilon_{\rm IR}$, $\ztrans$, $b$) with flat priors
alongside the standard cosmological parameters. The second,
\textit{EU} ($\Delta k = 0$), fixes all EU parameters at
the framework-derived values of Sec.~\ref{sec:params}.

\textit{Phase~1} runs each configuration with the IDE
perturbation source terms disabled, generating stable
proposal covariance matrices.
\textit{Phase~2} activates the source terms and uses the
Phase~1 covariance as the proposal, producing the baseline
posteriors.
\textit{Phase~3} extends the EU configuration with local
measurements:
\textit{EU+$H_0$} adds the SH0ES Cepheid distance-ladder
measurement~\cite{breuval2024};
\textit{EU+$H_0$+$S_8$} adds DES-Y3
$3 \times 2$pt cosmic shear via
Cocoa/CosmoLike~\cite{des2022,krause2021}, sampling 23
additional nuisance parameters (photometric redshifts,
shear calibration, intrinsic alignments, and galaxy bias).

All six runs share a baseline likelihood consisting of
Planck PR4 NPIPE CamSpec TTTEEE, low-$\ell$ TT and EE,
CMB lensing
$\phi\phi$~\cite{planck2020npipe}, DESI BAO
DR2~\cite{desi_dr2_2025}, and Pantheon+ SNe
Ia~\cite{scolnic2022}. Both the EU and $\Lambda$CDM
sample the same seven parameters: the six
$\Lambda$CDM cosmological parameters plus the Planck
calibration nuisance $A_{\rm planck}$.
Convergence is assessed via the Gelman--Rubin
statistic~\cite{gelman1992}; all six runs satisfy
$R-1 < 0.01$.

% ----------------------------------------------------------------
% V.B — Perturbation null test
% ----------------------------------------------------------------
\subsection{Perturbation null test}
\label{sec:pert_null_test}

The perturbation stability analysis of
Sec.~\ref{sec:stability} proved that the IDE source terms
cancel algebraically for $w = -1$ at the level of a single
CLASS evaluation. Phase~1 was designed to verify this
cancellation at the MCMC level: by running each
configuration with the source terms disabled, the
resulting posteriors provide a direct comparison against
the Phase~2 production runs.
All posterior shifts between Phase~1 and Phase~2 are
below $0.03\sigma$, consistent with Monte Carlo noise
for $N_{\rm eff} \simeq 33{,}000$ effective samples.
The Phase~1 and Phase~2 posteriors are statistically
indistinguishable; the full parameter-by-parameter
comparison is presented in
Table~\ref{tab:null_test} (Appendix~\ref{app:null_test}).
The four Phase~2 and Phase~3 posteriors therefore
constitute the complete analysis.

% [TABLE: Consolidated MCMC — ΛCDM + 4 EU runs]
\begin{table*}[t]
\caption{Posterior constraints from the four EU configurations,
compared with the $\Lambda$CDM best-fit obtained from the same
baseline datasets via minimizer. The EU column is the principal
result ($\Delta k = 0$, framework-derived parameters).
Uncertainties are $68\%$ credible intervals.}
\label{tab:mcmc}
\begin{tabular}{lccccc}
 & $\Lambda$CDM & EU-blind ($\Delta k{=}3$)
 & \textbf{EU} ($\Delta k{=}0$) & EU+$H_0$ & EU+$H_0$+$S_8$ \\
\hline
\multicolumn{6}{c}{\textit{Sampled parameters}} \\
$\omega_b$
  & $0.02238$ & $0.02231 \pm 0.00012$
  & $\mathbf{0.02218 \pm 0.00013}$ & $0.02221 \pm 0.00012$
  & $0.02218 \pm 0.00012$ \\
$\omega_{\rm cdm}$
  & $0.11933$ & $0.11778 \pm 0.00062$
  & $\mathbf{0.11926 \pm 0.00063}$ & $0.11913 \pm 0.00061$
  & $0.11944 \pm 0.00056$ \\
$100\theta_s$
  & $1.04097$ & $1.04191 \pm 0.00023$
  & $\mathbf{1.04180 \pm 0.00023}$ & $1.04181 \pm 0.00023$
  & $1.04182 \pm 0.00023$ \\
$\tau_{\rm reio}$
  & --- & $0.059 \pm 0.007$
  & $\mathbf{0.056 \pm 0.007}$ & $0.056 \pm 0.007$
  & $0.058 \pm 0.007$ \\
$\ln(10^{10}A_s)$
  & --- & $3.047 \pm 0.014$
  & $\mathbf{3.044 \pm 0.014}$ & $3.045 \pm 0.014$
  & $3.048 \pm 0.014$ \\
$n_s$
  & --- & $0.9676 \pm 0.0034$
  & $\mathbf{0.9637 \pm 0.0034}$ & $0.9640 \pm 0.0034$
  & $0.9641 \pm 0.0033$ \\[4pt]
\hline
\multicolumn{6}{c}{\textit{Derived parameters}} \\
$H_0^{\rm GKI}$ (km\,s$^{-1}$\,Mpc$^{-1}$)
  & $67.82$ & $68.13 \pm 0.28$
  & $\mathbf{68.89 \pm 0.28}$ & $68.96 \pm 0.27$
  & $68.83 \pm 0.25$ \\
$H_0^{\rm LKI}$ (km\,s$^{-1}$\,Mpc$^{-1}$)
  & --- & $71.90 \pm 0.29$
  & $\mathbf{72.71 \pm 0.30}$ & $72.78 \pm 0.28$
  & $72.64 \pm 0.26$ \\
$\sigma_8$
  & $0.803$ & $0.806 \pm 0.006$
  & $\mathbf{0.827 \pm 0.006}$ & $0.827 \pm 0.006$
  & $0.830 \pm 0.006$ \\
$S_8$
  & $0.812$ & $0.810 \pm 0.008$
  & $\mathbf{0.811 \pm 0.008}$ & $0.810 \pm 0.008$
  & $0.815 \pm 0.007$ \\
$\Omega_m$
  & $0.306$ & $0.303 \pm 0.004$
  & $\mathbf{0.288 \pm 0.003}$ & $0.288 \pm 0.003$
  & $0.289 \pm 0.003$ \\
$r_{\rm drag}$ (Mpc)
  & $147.86$ & $147.76 \pm 0.19$
  & $\mathbf{147.50 \pm 0.19}$ & $147.52 \pm 0.19$
  & --- \\
$f_{\rm cdm}(0)$
  & $1$ & $0.921 \pm 0.053$
  & $\mathbf{0.9558}$ & $0.9558$
  & $0.9558$ \\[4pt]
\hline
\multicolumn{6}{c}{\textit{Diagnostics}} \\
$\chi^2_{\rm total}$
  & $12{,}399.4$ & $12{,}394.4$
  & $\mathbf{12{,}396.3}$ & $12{,}396.5$
  & $13{,}026.6$ \\
$R-1$
  & --- & $0.003$ & $\mathbf{0.005}$ & $0.008$ & $0.006$ \\
$N_{\rm samples}$
  & --- & $63{,}887$ & $\mathbf{83{,}236}$ & $95{,}548$
  & $41{,}986$ \\
\end{tabular}
\end{table*}

% ----------------------------------------------------------------
% V.C — Posteriors
% ----------------------------------------------------------------
\subsection{EU parameters: free versus derived}
\label{sec:posteriors}

Table~\ref{tab:mcmc} presents the full posterior constraints.
The $\Lambda$CDM column reports the best-fit from a Cobaya
minimizer (BOBYQA, best of 4 starts) applied to the baseline
datasets.


The comparison between the EU-blind and EU columns of
Table~\ref{tab:mcmc} encodes the central statistical argument
of this analysis.

\textit{EU-blind $\to$ EU.}~The EU-blind run tests whether the data actively prefer evaporation.
The posterior $\varepsilon_{\rm IR} = 0.075 \pm 0.043$
excludes zero drain and is consistent with the
framework-derived value at $0.75\sigma$. However, the three
EU parameters are not individually constrained. Flat priors
on ($\varepsilon_{\rm IR}$, $z_{\rm trans}$, $b$) create a
three-dimensional degeneracy: geometrically distinct
combinations of transition redshift and coupling strength
produce nearly identical background expansion histories.
Current CMB, BAO, and SNe data constrain the integrated
drain fraction $f_{\rm cdm}(z{=}0)$ but lack the redshift
resolution to separate $\varepsilon_{\rm IR}$,
$z_{\rm trans}$, and $b$ individually --- all three
posteriors are prior-dominated (coefficient of variation
${>}\,47\%$; $f_{\rm cdm} = 0.921 \pm 0.053$).
The sampler explores a disproportionately large prior
volume of combinations that inflate the CDM drain beyond
its framework-derived value.
To preserve the acoustic peak structure under excess
evaporation, the MCMC compensates by lowering
$\omega_{\rm cdm}$ to $0.1178$ --- a $3.6\sigma$
displacement from the Planck $\Lambda$CDM inference ---
which suppresses $H_0^{\rm GKI}$ to
$68.13\;\text{km}\,\text{s}^{-1}\,\text{Mpc}^{-1}$.

Fixing the three EU parameters at their framework-derived
values eliminates this degeneracy. A single evaporation
curve --- a drain of exactly $4.42\%$ of CDM since
$z_{\rm trans} \approx 6$ --- replaces the cloud of
statistically degenerate alternatives. The primordial CDM
density recovers to $\omega_{\rm cdm} = 0.1193$
($1.2\sigma$ from Planck), and $H_0^{\rm GKI}$ rises by
$0.76\;\text{km}\,\text{s}^{-1}\,\text{Mpc}^{-1}$ to $68.89$. The best-fit $\chi^2$
increases by only $1.9$ for three fewer degrees of freedom;
the Akaike information criterion accordingly favors the
theory-fixed configuration
($\Delta\text{AIC} = -4.1$; see
Appendix~\ref{app:free_params}).



% ----------------------------------------------------------------
% V.D — Knockout tests
% ----------------------------------------------------------------
\subsection{Confrontation with local measurements}
\label{sec:knockout}

The progressive addition of local measurements tests
the rigidity of the EU posteriors.

\textit{Adding SH0ES (EU+$H_0$).}~The SH0ES Cepheid
distance-ladder measurement
$H_0 = 73.17 \pm 0.86\;\text{km}\,\text{s}^{-1}\,\text{Mpc}^{-1}$~\cite{breuval2024}
is incorporated via a custom likelihood that compares
against $H_0^{\rm LKI}$, the EU local observable
(Sec.~\ref{sec:lki}). All parameter shifts remain
$\leq 0.25\sigma$ (Table~\ref{tab:mcmc}), and the total
$\chi^2$ changes by only $0.2$: the CMB, BAO, and SNe
likelihoods are unperturbed by the addition of a local
measurement that the model had already accommodated. The
$S_8$ posterior shifts by $0.15\sigma$ ($0.811 \to 0.810$),
confirming that the $H_0$ and $S_8$ resolutions operate
through independent channels --- the CDM drain modifies the
geometric distance ladder without altering the amplitude of
structure growth. The SH0ES likelihood contributes no new
information: it is redundant for a model that predicts
$H_0^{\rm LKI} = 72.71\;\text{km}\,\text{s}^{-1}\,\text{Mpc}^{-1}$ from CMB, BAO,
and SNe alone.

\textit{Adding DES-Y3 (EU+$H_0$+$S_8$).}~This
configuration tests whether the EU posteriors
absorb weak-lensing data. The DES-Y3 $3 \times 2$pt cosmic
shear likelihood~\cite{des2022} introduces 23 additional
nuisance parameters (38 sampled dimensions total) and 462
data points. The maximum parameter shift relative to EU is
$0.50\sigma$ (in $S_8$), from $0.811$ to $0.815$; all other
parameters remain within $0.3\sigma$. The upward pull on
$S_8$ is expected: DES-Y3 directly constrains the
combination $\sigma_8 \sqrt{\Omega_m / 0.3}$, and the
joint fit rebalances the amplitude of structure growth
against the geometric expansion history. The $S_8$ tension
with the DES-Y3 published value decreases from $2.48\sigma$
($\Lambda$CDM) to $2.14\sigma$ (EU) --- a modest but
physical reduction driven by the lower $\Omega_m$ of the
drain scenario (Sec.~\ref{sec:drain_wl}). The maximum shift
between EU+$H_0$ and EU+$H_0$+$S_8$ is $0.49\sigma$ (in
$S_8$), confirming that the model remains stable under the
addition of an independent dataset probing a different
cosmological observable.

A caveat applies: the non-linear
matter power spectrum is computed via
HMCode~\cite{mead2021}, calibrated on $\Lambda$CDM
$N$-body simulations. Conservative scale cuts are applied
to restrict the analysis to the regime where EU
perturbations are algebraically exact
(Sec.~\ref{sec:pert_null_test}).

% ----------------------------------------------------------------
% V.C — UV→MCMC concordance
% ----------------------------------------------------------------
\subsection{Concordance with analytical predictions}
\label{sec:concordance}

The analytical predictions of Sec.~\ref{sec:gki_lki},
derived from the five framework parameters without
reference to observational data, are compared with the
MCMC posteriors in Table~\ref{tab:concordance}.


The maximum shift is $0.05\sigma$. The analytical framework
predicted $H_0^{\rm LKI} = 72.72\;\text{km}\,\text{s}^{-1}\,\text{Mpc}^{-1}$ before
the MCMC explored $83{,}236$ samples across three
independent datasets; the posterior mean converged to
$72.71 \pm 0.30\;\text{km}\,\text{s}^{-1}\,\text{Mpc}^{-1}$. This concordance
establishes the causal direction of the EU: QFT $\to$
framework parameters $\to$ CLASS $\to$ $H(z)$ $\to$ $H_0$.
The data do not inform the EU parameters --- they validate
that the framework-derived values are consistent with the
observable universe. The framework constrains the model more
tightly than the data alone --- the opposite of overfitting.

% [TABLE: UV vs MCMC concordance]
\begin{table}[pos=!htbp]
\caption{Concordance between the analytical predictions
(Sec.~\ref{sec:gki_lki}) and the MCMC posteriors (EU run).
No EU parameter is fitted to data.}
\label{tab:concordance}
\begin{tabular}{lccc}
Quantity & Analytic & MCMC (EU) & Shift \\
\hline
$H_0^{\rm GKI}$ (km\,s$^{-1}$\,Mpc$^{-1}$) & 68.90 & $68.89 \pm 0.28$ & $0.04\sigma$ \\
$H_0^{\rm LKI}$ (km\,s$^{-1}$\,Mpc$^{-1}$) & 72.72 & $72.71 \pm 0.30$ & $0.05\sigma$ \\
$f_{\rm cdm}(0)$ & 0.9558 & 0.9558 & exact \\
$\mathcal{I}_{\rm GKI}$ & 0.0679 & 0.0679 & exact \\
\end{tabular}
\end{table}


% ----------------------------------------------------------------
% V.E — Model selection
% ----------------------------------------------------------------
\subsection{Model selection}
\label{sec:model_selection}

With identical parameter counts ($k_{\rm EU} = k_{\Lambda\rm CDM}
= 7$), the information criteria~\cite{burnham_anderson2002}
reduce to pure goodness-of-fit:
\begin{equation}
\Delta\text{AIC} = \Delta\text{BIC}
  = \Delta\chi^2\,.
\label{eq:aic}
\end{equation}
This is the cleanest possible comparison: no Occam penalty,
no prior-volume ambiguity. Two dataset configurations are
compared, both using the $\Lambda$CDM best-fit obtained via
Cobaya minimizer (BOBYQA) on identical datasets. The
baseline comprises CMB, BAO, and SNe; the extended
configuration adds the SH0ES distance-ladder measurement
under symmetric conditions --- both models receive the KBC
void outflow correction of
Wu \& Huterer~\cite{wu_huterer2017}:
\begin{align}
\Delta\text{AIC}_{\rm baseline}
  &= -3.09 \quad (\text{positive evidence})\,,
  \label{eq:aic_baseline} \\
\Delta\text{AIC}_{\rm baseline+SH0ES}
  &= -6.88 \quad (\text{strong evidence})\,.
  \label{eq:aic_shoes}
\end{align}
The baseline improvement is driven by CMB
($\Delta\chi^2_{\rm CMB} = -4.2$) and BAO
($\Delta\chi^2_{\rm BAO} = -2.6$), partially offset by SNe
($\Delta\chi^2_{\rm SNe} = +3.7$). Adding SH0ES
consolidates the advantage: the EU predicts the correct
local $H_0$ from geometry alone, converting what is a
$5\sigma$ tension for $\Lambda$CDM into a
$0.5\sigma$ agreement.

The DES-Y3 data are not included in the model
selection. The published
$S_8 = 0.776$~\cite{des2022} was obtained under
$\Lambda$CDM using DES data alone, while the EU posterior
uses five combined datasets; a fair $\chi^2$ comparison
requires both models to be evaluated on identical data,
which was not performed for the DES-Y3 component.

The parameter shifts documented in
Table~\ref{tab:mcmc} --- primordial parameters unchanged,
acoustic densities rebalanced, sound horizon angle shifted
--- constitute a characteristic three-layer signature of the
CDM drain. The physical interpretation is developed in
Sec.~\ref{sec:drain}.



% ================================================================
% VI. DRAIN SIGNATURE
% ================================================================
\section{The Drain Signature in Cosmological Observables}
\label{sec:drain}

The MCMC posteriors of Sec.~\ref{sec:posteriors} are not
``Planck values.'' They are inferences conditioned on a
model. The six parameters reported by Planck~\cite{planck2018}
--- $\omega_b$, $\omega_{\rm cdm}$, $100\theta_s$,
$\tau_{\rm reio}$, $\ln(10^{10}A_s)$, $n_s$ --- were
extracted under the assumption that CDM is conserved.
Replacing that assumption with the EU drain modifies the
late-time expansion history, and the same CMB data respond
by readjusting the sampled parameters. The pattern of
these shifts is not arbitrary: it encodes the physics
of the CDM drain across three distinct layers.

% ----------------------------------------------------------------
% VI.A — Three-layer hierarchy
% ----------------------------------------------------------------
\subsection{Three-layer parameter hierarchy}
\label{sec:three_layers}

Table~\ref{tab:mcmc} reveals a characteristic hierarchy
when comparing the EU posteriors against the
$\Lambda$CDM best-fit.

\textit{Layer~1 --- Primordial ($< 0.3\sigma$).}~The
parameters $\ln(10^{10}A_s)$, $n_s$, and $\tau_{\rm reio}$
describe inflationary perturbations ($z \sim 10^{25}$) and
reionization ($z \sim 6$--$10$). The CDM drain operates at
$z < \ztrans \approx 6$ with a strict Heaviside cutoff
(Sec.~\ref{sec:params}). Physics above $\ztrans$ is
identical to $\Lambda$CDM, so the MCMC has no reason to
move these parameters. Their stability ($< 0.3\sigma$)
confirms that the EU preserves the primordial sector.

\textit{Layer~2 --- Acoustic.}~The baryon density
$\omega_b$ shifts by $-1.5\sigma$ relative to the
baseline $\Lambda$CDM fit (Table~\ref{tab:mcmc}), while
the primordial CDM density $\omega_{\rm cdm}$ remains
remarkably stable ($0.11926$ vs $0.11933$; shift of
$\sim 0.1\sigma$). The ratio
$\omega_b / \omega_{\rm cdm}$ controls the baryon
loading that governs the odd-even peak asymmetry in the
CMB power spectrum~\cite{hu_dodelson2002}; under the EU
expansion history, the MCMC rebalances $\omega_b$ to
preserve the peak structure measured by CamSpec.
The resulting $\omega_{\rm cdm} = 0.1193$ is the
primordial CDM density at recombination
($z \sim 1100$). The effective density today, after
$4.4\%$ evaporation, is
\begin{equation}
\omega_{\rm cdm}^{\rm eff}(0) =
  \omega_{\rm cdm}\, f_{\rm cdm}(0)
  = 0.1193 \cdot 0.9558 = 0.1140\,.
\label{eq:omega_eff}
\end{equation}
This $4.4\%$ reduction relative to the primordial
value is the origin of the reduced $\Omega_m$ and the
suppressed structure growth that alleviates the $S_8$
tension.

\textit{Layer~3 --- Geometric ($+3.6\sigma$).}~The
sound-horizon angle
$100\theta_s = r_s(z_*) / D_A(z_*)$ shifts from
$1.04097$ ($\Lambda$CDM) to $1.04180$ (EU). The CDM drain
transfers energy to the vacuum sector, accelerating the
late-time expansion and reducing $D_A(z_*)$. At fixed
sound horizon, the reduced angular diameter distance
requires a higher $H_0$ to preserve the observed CMB
peak positions --- yielding
$H_0^{\rm GKI} = 68.89\;\text{km}\,\text{s}^{-1}\,\text{Mpc}^{-1}$ instead of
$67.82$.

The hierarchy --- primordial frozen, densities rebalanced,
geometry shifted --- is the expected signature of a
late-time mechanism that modifies the expansion history
while leaving recombination physics intact.
% [TABLE: Mode A vs Mode B]
\begin{table}[h!]
\caption{CLASS-EU observables in two input cosmologies.
Mode~A: Planck $\Lambda$CDM best-fit parameters
($\omega_{\rm cdm} = 0.1200$, $\theta_s = 1.04092$).
Mode~B: MCMC posteriors
($\omega_{\rm cdm} = 0.1193$, $\theta_s = 1.04180$).
Both modes use identical EU parameters. The drain fraction
and perturbation cancellation are algebraic invariants.}
\label{tab:mode_ab}
\begin{tabular}{lccc}
Quantity & Mode A & Mode B & $\Delta$ \\
\hline
$H_0^{\rm GKI}$ (km\,s$^{-1}$\,Mpc$^{-1}$) & 68.56 & 68.89 & $+0.33$ \\
$\sigma_8$ & 0.829 & 0.827 & $-0.002$ \\
$S_8$ & 0.819 & 0.811 & $-0.008$ \\
$\Omega_m$ & 0.293 & 0.288 & $-0.005$ \\
$r_d$ (Mpc) & 147.06 & 147.42 & $+0.36$ \\
\hline
$f_{\rm cdm}(0)$ & 0.9558 & 0.9558 & $0$ \\
$\Delta\sigma_8$ (null test) & 0 & 0 & $0$ \\
\end{tabular}
\end{table}


% ----------------------------------------------------------------
% VI.B — Self-consistent observables
% ----------------------------------------------------------------
\subsection{Self-consistent observables}
\label{sec:mode_b}

The parameter readjustment has a concrete consequence for
all downstream observables. To isolate it, we compare the
fiducial CLASS-EU output of Sec.~\ref{sec:class_implementation}
--- computed with Planck $\Lambda$CDM input parameters
(Mode~A) --- against a second evaluation (Mode~B) that
replaces only the standard cosmological inputs
($\omega_b$, $\omega_{\rm cdm}$, $100\theta_s$) with the
MCMC posterior means of Table~\ref{tab:mcmc}. The five EU
parameters are identical in both cases. Each mode is a
single CLASS-EU evaluation, not an additional MCMC run.

The comparison exposes a diagnostic paradox:
the analytical prediction of Sec.~\ref{sec:gki_lki}
gives $H_0^{\rm GKI} = 68.90\;\text{km}\,\text{s}^{-1}\,\text{Mpc}^{-1}$; the MCMC
converges to $68.89 \pm 0.28$; but Mode~A yields $68.56$.

The resolution is that the Planck $\Lambda$CDM parameters
were calibrated for a universe without evaporation. Forcing
the EU physics into those inputs creates an inconsistent
hybrid --- the expansion history is modified, but the
cosmological densities have not relaxed. Mode~B produces
the self-consistent cosmology: $H_0 = 68.89$, in agreement
with the analytical prediction to $0.04\sigma$
(Table~\ref{tab:concordance}).

Table~\ref{tab:mode_ab} compares the two modes. The
$0.33\;\text{km}\,\text{s}^{-1}\,\text{Mpc}^{-1}$ offset in $H_0$ traces directly
to the acoustic-layer shift in $\omega_{\rm cdm}$ and
$\theta_s$. The GKI boost increases from $1.78\%$
(Mode~A) to $2.27\%$ (Mode~B) as the self-consistent
parameters allow the drain to operate at its analytically
predicted strength. All observables computed from Mode~B
--- $\sigma_8 = 0.827$, $S_8 = 0.811$,
$\Omega_m = 0.288$, $r_d = 147.42\;\text{Mpc}$ ---
constitute the definitive EU predictions.

Two quantities are invariant across modes. The drain
fraction $f_{\rm cdm}(0) = 0.9558$ depends only on the
EU parameters ($\varepsilon_{\rm IR}$, $\ztrans$, $b$),
which are identical in both configurations. The
perturbation cancellation $\Delta\sigma_8 = 0$
(Sec.~\ref{sec:stability}) is an algebraic property
of the IDE stress-energy tensor, independent of the
background cosmology. These invariants confirm that the
EU modification is fully specified by the framework
parameters --- the input cosmology determines only how
the standard $\Lambda$CDM parameters accommodate the
drain.

% ----------------------------------------------------------------
% VI.C — Bridge to IX
% ----------------------------------------------------------------
\subsection{Implications for structure growth}
\label{sec:drain_wl}

The drain signature propagates into the derived
observable $S_8 \equiv \sigma_8\sqrt{\Omega_m/0.3}$
through two competing effects. Relative to the Planck
2018 $\Lambda$CDM inference~\cite{planck2018}, the
perturbation amplitude $\sigma_8$ increases from $0.811$
to $0.827$ (EU) because the lower $\Omega_m$ reduces the
transfer-function suppression at late times. But
$\Omega_m$ itself decreases from $0.315$ to $0.288$,
and the square-root factor dominates: $S_8$ falls from
$0.831$ to $0.811$.

The reduced $S_8$ moves the EU prediction toward the
weak-lensing measurements. However, a direct comparison
between $S_8^{\rm EU}$ and the published survey values
is not valid. The parameter $S_8$ reported by DES-Y3,
KiDS, and HSC is not a model-independent observable: it
is inferred by fitting shear correlation functions
to a $\Lambda$CDM template with a fixed matter power
spectrum and a $\Lambda$CDM lensing kernel $W(\chi)$.
If the true cosmology is EU, the
$\Omega_m$-dependent kernel and the modified growth
history bias the inferred $S_8$. The magnitude of this
comparison artifact and its resolution are developed in
Sec.~\ref{sec:growth}.

% ================================================================
% VII. HUBBLE TENSION
% ================================================================
\section{Resolution of the Hubble Tension}
\label{sec:hubble}

The Hubble tension --- the $5.8\sigma$ discrepancy between
$H_0 = 67.36\;\text{km}\,\text{s}^{-1}\,\text{Mpc}^{-1}$ (Planck
$\Lambda$CDM~\cite{planck2018}) and
$H_0 = 73.17 \pm 0.86\;\text{km}\,\text{s}^{-1}\,\text{Mpc}^{-1}$
(SH0ES~\cite{breuval2024}) --- has resisted resolution by
late-time modifications to the Friedmann equation.
The five EU parameters derived in Sec.~\ref{sec:params}
predict $H_0^{\rm GKI} = 68.90\;\text{km}\,\text{s}^{-1}\,\text{Mpc}^{-1}$
(Sec.~\ref{sec:gki_lki}). Combined with the KBC void
outflow, the predicted local value is
$H_0^{\rm LKI} = 72.72\;\text{km}\,\text{s}^{-1}\,\text{Mpc}^{-1}$ ---
within $0.52\sigma$ of SH0ES, with $\Delta k = 0$.
The MCMC posteriors of Sec.~\ref{sec:posteriors}
confirm this prediction at $68.89 \pm 0.28$;
the comparisons that follow use the framework-derived
values throughout.

% ----------------------------------------------------------------
% VII.A — The geometric wall
% ----------------------------------------------------------------
\subsection{The geometric wall}
\label{sec:nogo}

Why has the Hubble tension resisted resolution? The answer is a
geometric constraint~\cite{schoneberg2022} that we
quantify via the profile likelihood.

We scan $H_0$ from $64$ to $76\;\text{km}\,\text{s}^{-1}\,\text{Mpc}^{-1}$
and, at each fixed $H_0$, minimize $\chi^2$ over
all remaining parameters (profiling over CMB~TT,
BAO~DR2, Pantheon+, and cosmic chronometers).
The resulting profile identifies the best achievable
fit as a function of $H_0$ for any model whose only
freedom is the Friedmann equation
(Fig.~\ref{fig:wall_escape}(a)). For $\Lambda$CDM,
the profile minimum lies at
$H_0^{\rm best} = 68.84$;
for the EU (GKI channel only), at $69.55$. As
$H_0$ is pushed toward $73.17$ (SH0ES), the
$\chi^2$ penalty grows quadratically:
\begin{equation}
\Delta\chi^2(H_0 = 73.17) = 158.41
\quad \Rightarrow \quad 12.59\sigma\,.
\label{eq:nogo_wall}
\end{equation}
This $12.59\sigma$ wall is model-independent in the
following sense: it applies to \textit{any} theory
that modifies only $H(z)$ while preserving the
pre-recombination physics. Early dark
energy~\cite{poulin2019_ede}, interacting dark
energy~\cite{divalentino2020},
$w_0 w_a$CDM~\cite{desi_w0wa}, and decaying dark
matter all face the same geometric constraint.
The wall exists because raising $H_0$ globally
shrinks both $r_s(z_*)$ and $D_A(z_*)$, and the
CMB peak positions pin their ratio to $\sim\!0.1\%$
precision.

\begin{figure}[pos=!htbp]
\includegraphics[width=\columnwidth]{fig1_wall_escape.pdf}
\caption{\textbf{(a)}~The geometric wall: $\Delta\chi^2$ profiles
for $\Lambda$CDM (solid) and the EU GKI (dashed). The SH0ES value
($H_0 = 73.17$) lies at $\Delta\chi^2 > 150$ ($12.59\sigma$).
\textbf{(b)}~Waterfall decomposition: Planck baseline ($67.36$),
GKI shift ($+1.53$), LKI void boost ($+3.82$), residual to
SH0ES ($0.52\sigma$).}
\label{fig:wall_escape}
\end{figure}

% ----------------------------------------------------------------
% VII.B — Two-scale resolution
% ----------------------------------------------------------------
\subsection{Two-scale resolution}
\label{sec:two_scale}

The EU respects the geometric wall rather than
attempting to breach it. The resolution separates
two physically distinct scales: a global kinematic
imprint (GKI) from the CDM drain, and a local
kinematic imprint (LKI) from the KBC void outflow
(Fig.~\ref{fig:wall_escape}(b)).

The GKI modifies the global Friedmann equation.
The CDM drain reduces $\Omega_{\rm cdm}$ at
$z < \ztrans$, and flatness requires $H_0$ to
increase to compensate. The framework predicts
$H_0^{\rm GKI} = 68.90$ (Sec.~\ref{sec:gki_lki});
the MCMC posteriors confirm
$H_0^{\rm GKI} = 68.89 \pm 0.28\;\text{km}\,\text{s}^{-1}\,\text{Mpc}^{-1}$
(Sec.~\ref{sec:posteriors}).
This shift is anchored by the CMB:
$\omega_{\rm cdm}$ is fixed at recombination, and the
drain fraction $f_{\rm cdm}(0) = 0.9558$ is determined
by the framework parameters. The GKI operates on the
entire observable universe and affects all cosmological
probes uniformly. At $H_0^{\rm GKI} = 68.90$, the
EU sits within the parabolic minimum of the profile
likelihood --- well below the geometric wall.

The LKI closes the remaining gap. The KBC void
($\delta_{\rm obs} = -0.46 \pm 0.06$, $R \approx
300\;\text{Mpc}$~\cite{keenan2013}) is an observed
underdensity surrounding the Milky Way. Within such a
void, bulk outflow biases distance-ladder measurements of
$H_0$~\cite{wu_huterer2017}. The magnitude of the bias
depends on the growth rate: $\delta H_0 / H_0 =
-(1/3)\,\delta_{\rm true}\,f(\Omega_m)$, where
$\delta_{\rm true} = -0.304$ after RSD
correction~\cite{haslbauer2020}. In the EU, the
growth rate $f(\Omega_m)$ is computed with the
self-consistent $\Omega_m = 0.288$, yielding
$\delta H_0 = +3.82\;\text{km}\,\text{s}^{-1}\,\text{Mpc}^{-1}$.
The void bias depends on $f(\Omega_m)$, which is
itself modified by the CDM drain: the same physics
that produces the GKI also determines the magnitude
of the LKI. The local boost is not an external
correction appended to the model --- it is a
consequence of the drain operating through structure
growth.

The two channels are additive and independent:
\begin{equation}
\begin{split}
H_0^{\rm LKI} &= H_0^{\rm GKI} + \delta H_0^{\rm void} \\
&= 68.90 + 3.82
= 72.72\;\text{km}\,\text{s}^{-1}\,\text{Mpc}^{-1}\,.
\end{split}
\label{eq:h0_lki}
\end{equation}
The GKI modifies the background cosmology; the LKI
modifies only the local frame of the observer. This
separation is rigid: the void outflow does not enter
the Friedmann equation and does not affect the CMB,
BAO, or any other global probe. It biases exclusively
the distance-ladder measurements that calibrate
supernovae within $\sim\!300\;\text{Mpc}$.

All local calibration methods --- Cepheids (SH0ES),
TRGB, JAGB, Miras --- share a common Rung~3
(supernovae in the Hubble flow) and therefore share
the void bias. The EU predicts that these methods
should cluster near $H_0^{\rm LKI} \approx 72.7$,
while CMB and inverse-distance-ladder measurements
should cluster near $H_0^{\rm GKI} \approx 68.9$. The
observed hierarchy (Table~\ref{tab:hierarchy}) is
consistent with this prediction
(Sec.~\ref{sec:hierarchy}).

% ----------------------------------------------------------------
% VII.C — Quantitative resolution
% ----------------------------------------------------------------
\subsection{Quantitative resolution}
\label{sec:quantitative_h0}

The EU reduces the Hubble tension from $5.8\sigma$
to $0.52\sigma$:
\begin{equation}
\frac{|H_0^{\rm LKI} - H_0^{\rm SH0ES}|}
     {\sigma_{\rm SH0ES}}
= \frac{|72.72 - 73.17|}{0.86}
= 0.52\sigma\,.
\label{eq:h0_tension}
\end{equation}
The numerator compares the framework prediction
($H_0^{\rm LKI} = 72.72$, Sec.~\ref{sec:gki_lki})
with the latest SH0ES calibration (Breuval et~al.
2024~\cite{breuval2024}). The EU prediction is
framework-derived; the only uncertainty enters from
the measurement. The MCMC posteriors confirm
$H_0^{\rm LKI} = 72.71 \pm 0.30$; including this
uncertainty in quadrature yields an identical
$0.51\sigma$.\footnote{Granting
$\Lambda$CDM the same KBC void correction
(Wu \& Huterer~\cite{wu_huterer2017}) yields
$H_0^{\Lambda{\rm CDM},\,{\rm local}} \approx 71.2\;\text{km}\,\text{s}^{-1}\,\text{Mpc}^{-1}$,
still in $2.3\sigma$ tension with SH0ES.
The EU advantage persists because the CDM drain
lowers $\Omega_m$ and shifts the growth rate
$f(\Omega_m)$, amplifying the GKI contribution
relative to $\Lambda$CDM.}

The EU+$H_0$ run (Phase~3, Table~\ref{tab:mcmc}) provides an independent
confirmation. Adding the SH0ES likelihood to the
baseline EU fit shifts $H_0^{\rm LKI}$ by
$0.25\sigma$, and the blind best-fit yields
$\chi^2(\text{SH0ES}) = 0.29$.
Both diagnostics indicate that the local
distance-ladder measurement is already accommodated
by the EU prediction without any pull on the
cosmological parameters. The SH0ES data are
redundant: the model predicted the correct local
$H_0$ before seeing it.

The LKI prediction is robust against the dominant
astrophysical uncertainty --- the KBC void contrast
$\delta$. The sensitivity band derived in
Sec.~\ref{sec:lki} spans $H_0^{\rm LKI} \in
[72.19,\, 73.26]$ at $1\sigma$, corresponding to a
worst-case tension of $1.14\sigma$ against SH0ES.

The cosmic age provides a consistency check:
$t_0 = 13.77\;\text{Gyr}$ (EU) versus
$13.80\;\text{Gyr}$ ($\Lambda$CDM). The difference
of $0.03\;\text{Gyr}$ ($0.2\%$) is negligible, and
both values exceed the age of the oldest globular
clusters ($t_{\rm GC} > 12.6\;\text{Gyr}$).

% ----------------------------------------------------------------
% VII.D — Distance probes
% ----------------------------------------------------------------
\subsection{Distance probes}
\label{sec:distance_probes}

The modified expansion history produces measurable
shifts in distance observables. We validate against
three independent datasets.

\textit{BAO --- DESI DR2.}~The 13-point DESI DR2
compilation~\cite{desi_dr2_2025} spans
$0.1 < z < 4.2$ and measures the transverse comoving
distance $D_M(z)/r_d$, the Hubble distance
$D_H(z)/r_d$, and the angle-averaged distance
$D_V(z)/r_d$. Both models are evaluated at their
respective best-fit parameters (Cobaya minimizer,
identical dataset suite). The EU achieves
$\chi^2_{\rm BAO} = 12.5$ compared with
$\chi^2_{\rm BAO} = 15.4$ for $\Lambda$CDM,
yielding
\begin{equation}
\Delta\chi^2_{\rm BAO} = -2.9
\quad (N_{\rm data} = 13)\,.
\label{eq:bao_dchi2}
\end{equation}
The physical origin is the modified late-time
expansion history. The EU sound horizon
$r_d = 147.50\;\text{Mpc}$ (versus $147.86$ for
$\Lambda$CDM) and the reduced matter fraction
$\Omega_m = 0.288$ (versus $0.306$) shift the
predicted distance ratios toward the DESI
measurements. The acoustic scale is determined at
recombination, where the EU and $\Lambda$CDM are
identical; the improvement comes from the
late-time expansion that converts $r_d$
into observed angles and redshifts.

\textit{Full-shape clustering --- DESI DR1.}~The
ShapeFit compressed parameters ($q_{\rm iso}$,
$q_{\rm AP}$, $df$, $dm$) measured by DESI
DR1~\cite{desi_dr1_shapefit} provide a complementary
test of the matter power spectrum shape. The EU
yields $\chi^2_{\rm ShapeFit} = 28.89$ versus
$32.39$ for $\Lambda$CDM
($\Delta\chi^2 = -3.50$, $n_{\rm dof} = 24$). The
growth suppression factor
$g \equiv \sigma_8^{\rm EU} / \sigma_8^{\rm fid}
= 0.987$ confirms that the CDM drain produces a
subtle but detectable modification of the growth
rate.

\textit{Type Ia supernovae and cosmic chronometers.}~The
Pantheon+ compilation~\cite{brout2022} contributes
$\chi^2 \approx 1408$ for both models --- the CDM drain
is too weak ($\varepsilon_{\rm IR} = 0.0426$) to produce
a measurable shift in distance moduli. The 32-point
cosmic-chronometer compilation~\cite{moresco2022}
yields $\Delta\chi^2 < 1$. Both datasets confirm that
the EU expansion history is conservative: it modifies
$H(z)$ at the percent level, preserving the
calibrated distance ladder.

% ================================================================
% VIII. N-BODY VALIDATION
% ================================================================
\section{N-body Validation}
\label{sec:nbody}

The analytical predictions and Boltzmann integration of the
preceding sections operate in the linear regime. Nonlinear
structure formation --- halo collapse, mergers, the nonlinear
power spectrum --- requires a full $N$-body simulation. We run
two parallel Gadget-4 simulations to validate that the CDM
drain is a pure background-level mass rescaling, decoupled
from the gravitational force computation.

% ----------------------------------------------------------------
% VIII.A — Simulation setup
% ----------------------------------------------------------------
\subsection{Simulation setup}
\label{sec:nbody_setup}

We use Gadget-4~\cite{gadget4} with a custom EU patch that
implements the CDM drain at the background level. The patch
reads a pre-computed table (\texttt{eu\_tables.h}) of
$f_{\rm cdm}(z)$ values generated from the analytical
formula of Sec.~\ref{sec:friedmann_eu} and rescales the CDM
particle mass at each time step. No modifications are
made to the gravitational force solver --- the
drain enters only through the mass assignment, not
through any fifth force or coupling in the Poisson equation.

Both simulations use $1024^3$ particles in a periodic box
of $500\;\text{Mpc}/h$, initialized at $z = 49$ with
second-order Lagrangian perturbation theory
(monofonIC~\cite{monofonic}). The cosmological parameters
are the baseline EU posteriors
(Table~\ref{tab:mcmc}):
$H_0 = 68.89\;\text{km}\,\text{s}^{-1}\,\text{Mpc}^{-1}$,
$\omega_{\rm cdm} = 0.1193$, $n_s = 0.965$,
$\ln(10^{10}A_s) = 3.044$. Power spectra are extracted at
15 snapshots from $z = 49$ to $z = 0$ using
Pylians3~\cite{pylians} with a $512^3$ CIC grid and
interlacing compensation.

Two configurations are run. The \textit{Production} run
uses the full TreePM gravitational solver; CDM particles
form halos, filaments, and voids as in a standard
cosmological simulation. The \textit{Null Test} run
disables self-gravity (\texttt{SELFGRAVITY = OFF}):
particles free-stream from their initial conditions,
generating no structure, while the CDM mass drain
continues to operate identically. From the Production
snapshots, a 2D interpolator $P(k, z)$ is constructed
(443 $k$-bins $\times$ 7 redshifts, $z = 0$ to $3$),
recovering exact $P(k)$ values at snapshot redshifts
to machine precision ($\sim\!10^{-14}\%$). This
interpolator feeds the downstream lensing analysis
(Sec.~\ref{sec:growth}) and the Stage~IV survey
predictions (Sec.~\ref{sec:predictions}).

% ----------------------------------------------------------------
% VIII.B — Null test
% ----------------------------------------------------------------
\subsection{Drain decoupling and the null test}
\label{sec:null_test}

The null test provides the most direct validation of the EU
implementation. With self-gravity disabled, the matter power
spectrum $P(k)$ remains flat shot noise across all 15
snapshots --- no halos, no filaments, no nonlinear
clustering. Yet the CDM mass evolves identically to the
Production run.

Table~\ref{tab:drain_nbody} presents the triple comparison
across all fifteen snapshots. The drain measured in the
Null Test, Production, and analytical formula agree to
$< 0.016\%$ at all epochs.

% [TABLE: Drain triple comparison]
\begin{table}[t]
\caption{CDM mass evolution: triple comparison between the
Null Test (self-gravity off), Production (full TreePM),
and the analytical formula from the EU framework
parameters. All masses are normalized to the $z = 49$
initial condition. The first three snapshots lie above
$\ztrans$ and show exactly zero drain.}
\label{tab:drain_nbody}
\begin{tabular}{lcccc}
$z$ & Null & Prod & Analytic
  & $|\Delta|_{\rm max}$ \\
\hline
$49.00$ & $1.0000$ & $1.0000$ & $1.0000$
  & $0\%$ \\
$6.99$  & $1.0000$ & $1.0000$ & $1.0000$
  & $0\%$ \\
$6.02$  & $1.0000$ & $1.0000$ & $1.0000$
  & $0\%$ \\
\hline
$5.50$  & $0.9759$ & $0.9758$ & $0.9758$
  & $0.016\%$ \\
$5.02$  & $0.9743$ & $0.9742$ & $0.9742$
  & $0.016\%$ \\
$4.50$  & $0.9724$ & $0.9723$ & $0.9723$
  & $0.016\%$ \\
$4.01$  & $0.9706$ & $0.9704$ & $0.9704$
  & $0.015\%$ \\
$3.51$  & $0.9685$ & $0.9684$ & $0.9684$
  & $0.014\%$ \\
$2.99$  & $0.9664$ & $0.9663$ & $0.9663$
  & $0.013\%$ \\
$2.50$  & $0.9644$ & $0.9643$ & $0.9643$
  & $0.011\%$ \\
$2.00$  & $0.9624$ & $0.9623$ & $0.9623$
  & $0.009\%$ \\
$1.50$  & $0.9604$ & $0.9603$ & $0.9603$
  & $0.007\%$ \\
$1.00$  & $0.9586$ & $0.9586$ & $0.9585$
  & $0.005\%$ \\
$0.50$  & $0.9570$ & $0.9570$ & $0.9570$
  & $0.003\%$ \\
$0.00$  & $0.9558$ & $0.9558$ & $0.9558$
  & $0.002\%$ \\
\end{tabular}
\end{table}

Three results follow from this comparison
(Fig.~\ref{fig:anatomy}(b)).

\textit{(i)~Gravitational decoupling.}~The drain is
identical whether or not particles interact
gravitationally. Dense clusters and empty voids
evaporate at the same fractional rate. The CDM drain
is a background-level thermodynamic process ---
it rescales the mass budget without modifying the
Poisson equation or introducing any fifth force.
This confirms at the nonlinear level the algebraic
cancellation proven analytically in
Sec.~\ref{sec:stability}.

\textit{(ii)~UV concordance.}~The total drain at
$z = 0$ is $4.42\%$ in both simulations, matching
the analytical prediction from the five
framework-derived parameters
(Sec.~\ref{sec:params}). The EFT prediction is
confirmed by $10^9$ particles to an accuracy of
$0.003\%$.

\textit{(iii)~Heaviside cutoff.}~Above $\ztrans$,
the mass spread across snapshots is exactly zero:
$f_{\rm cdm}(z{=}49) = f_{\rm cdm}(z{=}6.99) =
f_{\rm cdm}(z{=}6.02)$. The first mass drop occurs
at $z = 5.50$ ($2.4\%$ below the initial
condition), confirming that the Heaviside cutoff
in Eq.~(\ref{eq:eps}) is preserved exactly in
the nonlinear regime. No drain leaks into the
pre-recombination universe.

\begin{figure}[pos=!htbp]
\includegraphics[width=\columnwidth]{fig2_anatomy.pdf}
\caption{Anatomy of the evaporating universe. \textbf{(a)}~The
coupling function $\varepsilon(z)$, with logistic activation
at $\ztrans \approx 5.99$ and IR plateau
$\epsIR = 4.26\%$. \textbf{(b)}~CDM survival fraction
$f_{\rm cdm}(z)$: analytical solution (solid),
$N$-body snapshots (circles), and null-test
particles (diamonds).}
\label{fig:anatomy}
\end{figure}

% ----------------------------------------------------------------
% VIII.C — Power spectrum and growth
% ----------------------------------------------------------------
\subsection{Power spectrum suppression and anemic halos}
\label{sec:pk_suppression}

The suppression ratio
$S(k) \equiv P_{\rm EU}(k) / P_{\Lambda{\rm CDM}}(k)$
at $z = 0$ reveals three physically distinct regimes
(Fig.~\ref{fig:power_spectrum}), where
$P_{\Lambda{\rm CDM}}$ is computed from CLASS
with HMCode nonlinear corrections~\cite{mead2021}
using the same cosmological parameters.

\textit{Linear scales ($k < 0.1\;h/\text{Mpc}$).}~The
suppression ratio plateaus at $S \approx 1.065$,
consistent with the linear Conservative Upper Bound (CUB) expectation
$(\sigma_8^{\rm EU} / \sigma_8^{\Lambda\rm CDM})^2
= 1.038$ given the cosmic variance of the
$500\;\text{Mpc}/h$ simulation volume. The EU has
more linear power than $\Lambda$CDM because the
reduced $\Omega_m$ requires a higher $\sigma_8$ to
fit the same CMB acoustic peaks. This is the same
mechanism that raises $\sigma_8$ analytically in
Sec.~\ref{sec:drain_wl}.

\textit{Quasi-linear scales
($0.1 < k < 1.0\;h/\text{Mpc}$).}~$S(k)$ dips
below the CUB line. This is the signature of
\textit{anemic halos}: CDM evaporation reduces
the mass of collapsed objects, producing shallower
potential wells and less concentrated density
profiles. The power spectrum suppression at these
scales drives the $S_8$ reduction that alleviates
the growth tension (Sec.~\ref{sec:growth}).

\textit{Deep nonlinear scales
($k > 1.0\;h/\text{Mpc}$).}~$S(k)$ rises sharply
above the CUB. This is a known artifact of
dark-matter-only simulations: the $\Lambda$CDM
reference (HMCode) includes baryonic feedback
that suppresses halo cores, while the Gadget-4
run has no such feedback. The rise does not
reflect EU physics; it validates the decision to
anchor the small-scale power spectrum to HMCode
in the downstream $S_8$ analysis
(Sec.~\ref{sec:growth}).

\begin{figure}[pos=!htbp]
\includegraphics[width=\columnwidth]{fig3_power_spectrum.pdf}
\caption{Power spectrum suppression ratio
$S(k) \equiv P_{\rm EU}(k) / P_{\Lambda{\rm CDM}}(k)$ at $z = 0$.
Three regimes: linear plateau at $S \approx 1.065$
(large scales), anemic-halo dip (quasi-linear,
driving the $S_8$ resolution), and rise at
$k > 1\;h/\text{Mpc}$ from the absence of baryonic
feedback. Dashed line: CUB prediction.}
\label{fig:power_spectrum}
\end{figure}

The three-regime structure persists across redshifts.
The linear plateau grows monotonically from
$S = 1.044$ at $z = 2$ through $1.052$ ($z = 1$) and
$1.057$ ($z = 0.5$) to $1.065$ ($z = 0$), tracking
the accumulating CDM drain (Table~\ref{tab:drain_nbody}).
The quasi-linear dip deepens in parallel: halos at
$z = 0$ have lost $4.4\%$ of their initial CDM mass,
versus $3.8\%$ at $z = 2$, producing progressively
more anemic density profiles. This redshift dependence
is a prediction unique to the EU --- in $\Lambda$CDM,
the ratio $P_{\rm model}/P_{\rm ref}$ is unity at
all epochs by construction.

The linear growth factor confirms the suppression.
Extracting $D(z)$ from the $N$-body power spectrum
at $k_{\rm ref} = 0.05\;h/\text{Mpc}$ and comparing
with the analytical growth ODE yields
\begin{equation}
g \equiv \frac{D_{\rm EU}(0)}{D_{\Lambda{\rm CDM}}(0)}
= 0.992\,,
\label{eq:growth_suppression}
\end{equation}
corresponding to a $0.8\%$ suppression relative to
$\Lambda$CDM. The $N$-body extraction agrees with the
analytical ODE to $\Delta g = 0.0004$, confirming that
the linear growth regime is faithfully reproduced by
the simulation.

% ================================================================
% IX. GROWTH TENSION
% ================================================================
\section{Resolution of the Growth Tension}
\label{sec:growth}

The parameter $S_8 \equiv \sigma_8\sqrt{\Omega_m/0.3}$ reported
by weak-lensing surveys is not a model-independent observable.
It is extracted by fitting shear correlation functions to
$\Lambda$CDM templates --- using a $\Lambda$CDM lensing
kernel, a $\Lambda$CDM comoving distance, and a $\Lambda$CDM
matter power spectrum. If the true cosmology is EU, every
link in this inference chain carries a systematic bias.

The EU resolves the growth tension through two complementary
arguments. The EU+$H_0$+$S_8$ run (Sec.~\ref{sec:mcmc_confirmation})
proves numerically that DES-Y3 data are consistent with
$S_8 = 0.815$ when analyzed under EU cosmology. The forward
analysis explains the mechanism: geometric kernel
bias of $-5\%$ to $-7\%$ systematically depresses the
$S_8$ inferred by $\Lambda$CDM pipelines.

% ----------------------------------------------------------------
% IX.A — The lensing kernel bias
% ----------------------------------------------------------------
\subsection{The lensing kernel bias}
\label{sec:kernel_bias}

Weak-lensing surveys measure the angular cross-spectrum of
cosmic shear between tomographic bins $i$ and $j$. Under the
Limber approximation~\cite{limber1953},
\begin{equation}
C_\ell^{ij} = \int_0^{\chi_H}
  \frac{W_i(\chi)\,W_j(\chi)}{\chi^2}\;
  P\!\left(\frac{\ell + 1/2}{\chi},\, z(\chi)\right)
  d\chi\,,
\label{eq:limber}
\end{equation}
where $P(k,z)$ is the matter power spectrum and $W_i(\chi)$
is the lensing efficiency kernel for bin $i$:
\begin{equation}
W_i(\chi) = \frac{3\,\Omega_m\,H_0^2}{2\,c^2}\;
  \frac{\chi}{a(\chi)}\;
  \int_\chi^{\chi_H}
  n_i(\chi')\,\frac{\chi' - \chi}{\chi'}\;d\chi'\,.
\label{eq:kernel}
\end{equation}
Three quantities in Eq.~(\ref{eq:kernel}) depend on the
assumed cosmology: the matter density $\Omega_m$, the
comoving distance $\chi(z)$, and the source distribution
$n_i(\chi)$ mapped from redshift space. The EU modifications
--- $\Omega_m^{\rm EU} = 0.288$ (vs $0.315$) and
$\chi_{\rm EU}(z) \neq \chi_{\Lambda{\rm CDM}}(z)$ due to
the altered expansion history --- shift the kernel peaks,
change the overlap between tomographic bins, and modify
the Limber mapping $k = (\ell + 1/2)/\chi$.

A $\Lambda$CDM pipeline analyzing shear data from an EU
universe uses the wrong kernel in every integral. To
quantify the resulting bias, we define the
\textit{pipeline-inferred} value
\begin{equation}
\hat{S}_8 = S_8^{\rm EU}\,
  \sqrt{\left\langle A_{ij} \right\rangle_{\rm Fisher}}\,,
\label{eq:shat}
\end{equation}
where $A_{ij} = C_\ell^{ij}(\text{EU}) /
C_\ell^{ij}(\Lambda\text{CDM})$ is the amplitude ratio
of the true (EU) to template ($\Lambda$CDM) angular spectra,
and the Fisher average weights each bin pair by its
signal-to-noise ratio. The ratio $\hat{S}_8 / S_8^{\rm EU}$
is the kernel bias.

% ----------------------------------------------------------------
% IX.B — Resolution by survey
% ----------------------------------------------------------------
\subsection{Quantitative resolution}
\label{sec:s8_resolution}

We compute $\hat{S}_8$ for four independent Stage-III
surveys using our forward-modeling pipeline:
mock shear spectra are generated with the EU matter power
spectrum from the $N$-body simulation
(Sec.~\ref{sec:pk_suppression}) and EU lensing kernels,
then reanalyzed with $\Lambda$CDM kernels via
Fisher-weighted Knox variance~\cite{knox1995}.
Each survey uses its official source redshift distributions,
except HSC-Y3, for which we adopt a Smail
parametrization matched to the published
tomographic bins.

Table~\ref{tab:s8} and Fig.~\ref{fig:s8_forest} present the results. All four surveys
converge to $\hat{S}_8 \approx 0.794 \pm 0.002$, with
inter-survey scatter of $\Delta\hat{S}_8 = 0.004$ ---
smaller than any individual statistical uncertainty.

% [TABLE: S₈ tensions by survey]
\begin{table}[t]
\caption{Resolution of the growth tension across four
independent Stage-III weak-lensing surveys. $S_8^{\rm pub}$
is the published value (inferred under $\Lambda$CDM);
$\hat{S}_8$ is the pipeline-inferred value in an EU
universe (forward model). All kernel biases are
geometric and require no additional parameters.}
\label{tab:s8}
\begin{tabular}{lccccc}
Survey & $S_8^{\rm pub}$ & $\hat{S}_8$ &
  $\sigma_{\Lambda\rm CDM}$ &
  $\sigma_{\rm EU}$ & Bias \\
\hline
DES-Y3        & $0.776$ & $0.794$ &
  ${\sim}\,2.6\sigma$ & $\mathbf{-1.03\sigma}$ & $-5.6\%$ \\
KiDS-1000~\cite{heymans2021}
  & $0.759$ & $0.792$ &
  $2.17\sigma$ & $\mathbf{-1.38\sigma}$ & $-6.6\%$ \\
KiDS-Legacy~\cite{wright2025}
  & $0.815$ & $0.793$ &
  $0.89\sigma$ & $\mathbf{+1.15\sigma}$ & $-6.2\%$ \\
HSC-Y3~\cite{li2023_hsc}
  & $0.776$ & $0.796$ &
  $1.10\sigma$ & $\mathbf{-0.63\sigma}$ & $-5.4\%$ \\
\end{tabular}
\end{table}

The convergence across surveys with different source
distributions, sky coverage, and shear calibration confirms
that the kernel bias is a universal geometric effect of the
EU expansion history, not a survey-specific systematic.
The pattern is particularly instructive: DES-Y3 and
KiDS-Legacy bracket the EU expectation from opposite
sides ($-1.03\sigma$ and $+1.15\sigma$, respectively),
while KiDS-1000 and HSC-Y3 scatter below it. Under
$\Lambda$CDM, these surveys present an incoherent picture
--- DES-Y3 reports a ${\sim}\,2.6\sigma$ deficit while
KiDS-Legacy finds sub-$\sigma$ agreement. Under the EU,
all four measurements reduce to standard
${\sim}\,1\sigma$ statistical scatter around a common
geometric expectation.

The KiDS-Legacy result merits specific attention. Wright
et~al.~\cite{wright2025} replaced the direct calibration
method of KiDS-1000 with Self-Organizing Map (SOM)
photometric redshifts, calibrated against $126{,}085$
spectroscopic sources, and extended the tomographic
analysis to six bins ($z \leq 1.2$). While this
improved $n(z)$ calibration raised the published $S_8$
from $0.759$ to $0.815$, the geometric kernel bias
persists: the SOM calibrates $n(z)$ from spectral lines
(a purely kinematic measurement), but the Limber
projection in CosmoSIS continues to assume
$\chi_{\Lambda{\rm CDM}}(z)$.

The bias depends on the source redshift distribution.
Table~\ref{tab:kernel_bias} presents the per-bin kernel
bias for DES-Y3 and KiDS-Legacy, both with official
$n(z)$ distributions. For DES-Y3, the bias peaks at
$\bar{z} = 0.70$ ($-6.2\%$), where the divergence between
$\chi_{\rm EU}(z)$ and $\chi_{\Lambda{\rm CDM}}(z)$
maximizes. For KiDS-Legacy, the bias decreases
monotonically from $-7.1\%$ (BIN1, $\bar{z} \approx 0.34$)
to $-4.9\%$ (BIN6, $\bar{z} \approx 1.22$), reflecting
the reduced CDM drain fraction at higher redshifts.
This redshift structure is a prediction: Euclid, with 10
tomographic bins extending to $z \approx 2$, will map
this curve with percent-level precision
(Sec.~\ref{sec:euclid}).

\begin{table}[t]
\caption{Per-bin kernel bias for the two Stage-III surveys
with official source redshift distributions. The bias
$\mathcal{B}_i \equiv (W_i^{\rm EU} - W_i^{\Lambda{\rm CDM}})
/ W_i^{\Lambda{\rm CDM}}$ is the fractional deficit of the
EU lensing kernel relative to $\Lambda$CDM in each
tomographic bin.}
\label{tab:kernel_bias}
\begin{tabular}{lccc}
Survey & Bin & $\bar{z}$ range & Bias (\%) \\
\hline
DES-Y3 & 1 & $0.0$--$0.36$ & $-5.1$ \\
        & 2 & $0.36$--$0.63$ & $-5.7$ \\
        & 3 & $0.63$--$0.87$ & $-6.2$ \\
        & 4 & $0.87$--$1.30$ & $-5.3$ \\
\hline
KiDS-Legacy & 1 & $0.1$--$0.3$ & $-7.1$ \\
            & 2 & $0.3$--$0.5$ & $-7.1$ \\
            & 3 & $0.5$--$0.7$ & $-6.7$ \\
            & 4 & $0.7$--$0.9$ & $-6.1$ \\
            & 5 & $0.9$--$1.1$ & $-5.6$ \\
            & 6 & $1.1$--$1.5$ & $-4.9$ \\
\end{tabular}
\end{table}

The total shift from $S_8^{\Lambda\rm CDM} = 0.832$ to
$\hat{S}_8 = 0.794$ decomposes into two physical
contributions: the geometric kernel bias
($\Delta S_8 = -0.028$, $\sim\!75\%$) and the anemic halo
suppression from the quasi-linear $P(k)$ dip
($\Delta S_8 = -0.010$, $\sim\!25\%$;
Sec.~\ref{sec:pk_suppression}). The kernel bias dominates
because the $C_\ell$ integrand
(Eq.~\ref{eq:limber}) depends quadratically on $W(\chi)$:
a $6\%$ kernel deficit propagates as a $\sim\!12\%$
suppression in $C_\ell$, which the pipeline interprets
as a lower $\sigma_8$.

The $S_8$ ``ladder'' summarizes the chain of inferences:
\begin{align}
S_8^{\Lambda\rm CDM} &= 0.832
  \quad\text{(Planck, $\Lambda$CDM)} \notag \\
S_8^{\rm EU} &= 0.811
  \quad\text{(EU prediction, no WL data)} \notag \\
\hat{S}_8 &= 0.794
  \quad\text{(EU $\to$ $\Lambda$CDM pipeline)} \notag \\
S_8^{\rm pub} &= 0.776
  \quad\text{(DES-Y3 published)}\,.
\label{eq:s8_ladder}
\end{align}
The $\Lambda$CDM tension (typically reported at
${\sim}\,2.6\sigma$) arises from comparing two numbers
computed in different cosmologies. The EU tension
(${\sim}\,1\sigma$) compares like with like.


\begin{figure}[pos=!htbp]
\includegraphics[width=\columnwidth]{fig4_s8_forest.pdf}
\caption{Resolution of the growth tension across four independent
Stage-III weak-lensing surveys. For each survey, the published
$S_8^{\rm pub}$ (extracted under $\Lambda$CDM assumptions) is
compared with the EU forward-model prediction $\hat{S}_8$. The
kernel bias shifts all surveys toward concordance with the CMB,
reducing the average tension from $\sim\!2.6\sigma$ to
$\sim\!1\sigma$.}
\label{fig:s8_forest}
\end{figure}


A parallel pattern emerges in the growth data. Under
$\Lambda$CDM with constant $\Omega_m$, the CMB yields
$S_8 = 0.832 \pm 0.013$~\cite{planck2018} while
DES-Y3 cosmic shear reports
$S_8 = 0.776 \pm 0.017$~\cite{des2022} --- a
${\sim}\,2.6\sigma$ discrepancy between the early and
late universe. The CMB measurement reflects the
primordial matter density at $z \sim 1100$; the
weak-lensing measurement is sensitive to the
integrated growth rate at $z < 1$. When these
datasets are jointly analyzed under $\Lambda$CDM
(together with BAO, supernovae, and CMB lensing),
the single-$\Omega_m$ compromise converges to
$S_8 = 0.812 \pm 0.008$~\cite{des2022}.
The arithmetic mean of the two inputs is
$S_8 \approx 0.804$; the joint constraint lies
$1\sigma$ above this because the CMB carries an order
of magnitude more Fisher information than cosmic shear
at current survey depth.

The EU predicts $S_8 = 0.811$
(Sec.~\ref{sec:s8_resolution}), within $0.001$ of
this $\Lambda$CDM joint constraint. The agreement has
a structural explanation. In the EU, $\Omega_m(z)$
evolves: high at recombination (consistent with the
CMB), $4.4\%$ lower at $z = 0$ (consistent with
weak lensing). A $\Lambda$CDM analysis, constrained
to a single constant $\Omega_m$, has no choice but
to find the value that best approximates this
dynamical trajectory. The result is a compromise
pulled toward the high-weight CMB measurement but
dragged below it by the late-universe data --- landing
on the EU prediction.

% ----------------------------------------------------------------
% IX.C — Robustness
% ----------------------------------------------------------------
\subsection{Robustness tests}
\label{sec:s8_robustness}

\textit{UV extrapolation.}~The Gadget-4 simulation has
$k_{\rm Nyquist} \approx 4.3\;h/\text{Mpc}$. Beyond this
scale, we perform an A/B test replacing the default
power-law extrapolation with HMCode2020
anchoring~\cite{mead2021}. The maximum impact on
$\hat{S}_8$ is $0.074\sigma$ (DES-Y3); all four
surveys satisfy the criterion $\Delta\hat{S}_8/\sigma
< 0.2$. Stage-III surveys are insensitive to the
deep nonlinear power spectrum.

\textit{Intrinsic alignments.}~Adding the NLA
model~\cite{bridle2007} with amplitude $A_{\rm IA} = 1.0$
(DES-Y3 fiducial) shifts $\hat{S}_8$ upward by $+0.005$.
The tension with DES-Y3 increases from $1.03\sigma$ to
$1.33\sigma$ --- still well within concordance. Ignoring
intrinsic alignments is the conservative choice.

\textit{Residual gap.}~The forward analysis captures
$\sim\!70\%$ of the total $S_8$ shift ($0.832 \to 0.794$;
$\Delta = -0.038$ of a total $-0.056$). The remaining
gap ($0.794 \to 0.776$) arises from pipeline effects not
modeled in our Fisher analysis: baryonic feedback,
photo-$z$ uncertainties, shear calibration, and
nonlinear galaxy bias. All of these are included in the
full CosmoLike likelihood used by the EU+$H_0$+$S_8$ run, which
independently confirms $S_8 = 0.815 \pm 0.007$
(Sec.~\ref{sec:mcmc_confirmation}). The forward
analysis is a diagnostic tool that explains the
mechanism; the full MCMC is the definitive numerical
verdict.

% ----------------------------------------------------------------
% IX.D — MCMC confirmation
% ----------------------------------------------------------------
\subsection{MCMC confirmation}
\label{sec:mcmc_confirmation}

The EU+$H_0$+$S_8$ run adds the DES-Y3 $3 \times 2$pt cosmic shear
likelihood~\cite{des2022} (462 data points, analyzed via
Cocoa/CosmoLike~\cite{krause2021}) to the baseline
CMB+BAO+SNe+SH0ES dataset. The sampling space expands
to 38 dimensions: 7 cosmological parameters,
8 instrumental nuisance parameters ($A_{\rm planck}$
and Planck calibration), and 23 DES nuisance parameters
(shear calibration $m_i$,
photo-$z$ shifts $\Delta z_i$, intrinsic alignment
amplitudes $A_{\rm IA}$, galaxy bias $b_i$, and
baryonic feedback).

The posterior converges at $R-1 < 0.006$
(per-parameter Gelman--Rubin diagnostic) with
$41{,}986$ samples across eight chains
(Table~\ref{tab:mcmc}).
The key derived parameters are
\begin{align}
S_8^{\rm full} &= 0.815 \pm 0.007
  \quad (0.50\sigma\;\text{from baseline})\,,
  \notag \\
H_0^{\rm GKI} &= 68.83 \pm 0.25
  \quad (0.23\sigma\;\text{from baseline})\,,
  \notag \\
\Omega_m &= 0.289 \pm 0.003
  \quad (0.27\sigma\;\text{from baseline})\,.
\label{eq:d2_params}
\end{align}
Adding 462 weak-lensing data points shifts $S_8$ by
$0.50\sigma$ --- from $0.811$ (CMB+BAO+SNe alone) to
$0.815$ (including DES-Y3). The model absorbs the
full DES-Y3 $3 \times 2$pt dataset without stress.

This result resolves both caveats of the forward
analysis: the simplified background
integration and the diagonal Fisher approximation. The
full CosmoLike pipeline includes baryonic feedback,
photo-$z$ marginalization, shear calibration, and the
complete off-diagonal covariance matrix.

The EU growth tension is best quantified through the
three hierarchical $S_8$ values that the framework
produces. The first is the CLASS-EU standalone
prediction (Mode~B, Sec.~\ref{sec:mode_b}):
the five EFT-derived parameters are fed into the
Boltzmann solver with the baseline cosmological posteriors,
and $S_8$ is computed without any weak-lensing data
entering the calculation. The result is $S_8 = 0.811$.
The second is the baseline MCMC posterior mean,
$S_8 = 0.811 \pm 0.008$, obtained by marginalizing
over all six $\Lambda$CDM parameters while fitting
CMB, BAO, and supernovae simultaneously. The third is
the full posterior, $S_8 = 0.815 \pm 0.007$, which
adds 462 DES-Y3 data points and 23 nuisance parameters
to the fit.

These three values --- 0.811, 0.811, 0.815 --- span a
range of $0.004$, smaller than the statistical
uncertainty of any single determination. The analytical
prediction, the Phase~2 EU posterior, and the Phase~3
EU+$H_0$+$S_8$ posterior converge to the same $S_8$. The
shift from the blind prediction to the full fit,
\begin{equation}
\frac{|S_8^{\rm full} - S_8^{\rm Mode\,B}|}
     {\sigma_{\rm full}}
= \frac{|0.8151 - 0.8112|}{0.0067} = 0.58\sigma\,,
\label{eq:s8_tension_eu}
\end{equation}
is the true EU growth tension: a model predicted
$S_8 = 0.811$ without seeing any shear data, and the
same shear data, analyzed within that model, returned
$S_8 = 0.815$.

The contrast with $\Lambda$CDM is stark. Under
$\Lambda$CDM, the CMB predicts $S_8 = 0.832$, while
DES-Y3 measures $S_8 = 0.776 \pm 0.017$ --- a
${\sim}\,2.6\sigma$ internal tension with no mechanism
to bridge the gap. Under the EU, the framework predicts
$S_8 = 0.811$ before seeing any shear data; adding the
full DES-Y3 likelihood to the fit shifts the posterior
to $S_8 = 0.815 \pm 0.007$ --- a $0.58\sigma$
agreement.

For $\Lambda$CDM, the growth tension is genuine: the
model predicts $S_8 = 0.832$ and most Stage-III
surveys report $S_8 \lesssim 0.78$, with no internal
mechanism to bridge the gap. For the EU, the kernel
bias accounts for the dominant share of this apparent
deficit ($\hat{S}_8 \approx 0.794$;
Table~\ref{tab:s8}). The convergence of four
independent surveys to the same suppressed
$\hat{S}_8$ band is not a coincidence requiring
explanation --- it is a geometric consequence of
analyzing an evaporating universe with a
non-evaporating template.

% ----------------------------------------------------------------
% IX.E — Pipeline bias generalization
% ----------------------------------------------------------------
\subsection{Pipeline bias generalization}
\label{sec:pipeline_bias}

The forensic pipeline dissection performed above for $S_8$
revealed a \emph{model-intrinsic} bias: the lensing kernel
$W(\chi)$ carries $\Lambda$CDM-calibrated parameters that
systematically deflate the inferred $\hat{S}_8$.
This bias is monotonically redshift-dependent
(Table~\ref{tab:kernel_bias}), peaking at low $z$ where the
cumulative CDM drain is largest. Whether this redshift
structure contributes to the photo-$z$ calibration residuals
reported at low redshift by Stage-III surveys merits
independent investigation.

An analogous bias affects the Hubble constant, but its origin
is \emph{structural} rather than model-intrinsic --- the
SH0ES calibrators reside within the KBC
void~\cite{keenan2013}, whose gravitational outflow
contaminates $M_B$ and inflates the inferred $H_0$.
Because this cosmographic bias depends on the local void
profile rather than on the cosmological model, its forensic
decomposition requires independent probes (peculiar
velocities, BAO-anchored joint analysis) that fall outside
the scope of the present work.
The complete dissection will be presented in Paper~II.

% ================================================================
% X. FALSIFIABLE PREDICTIONS
% ================================================================
\section{Falsifiable Predictions}
\label{sec:predictions}

The EU framework, with its five derived parameters, generates
quantitative predictions for three independent Stage-IV
programs. Each prediction follows from the baseline EU posteriors
and the $N$-body calibration of the preceding sections.
No parameter is adjusted to the test data.

% ----------------------------------------------------------------
% X.A — Euclid DR1
% ----------------------------------------------------------------
\subsection{Euclid DR1 weak lensing and spectroscopy}
\label{sec:euclid}

\textit{Anomalous $C_\ell$ pattern.}~The lensing kernel bias
identified in Sec.~\ref{sec:kernel_bias} predicts a specific
suppression of the angular power spectrum measurable by
Euclid~\cite{euclid2024}. Across the 10 tomographic bins of the
Euclid photometric survey ($0.001 < z < 2.5$), the lensing
kernel $W(\chi)$ is suppressed by $-8.1\%$ (Bin~1,
$\bar{z} = 0.30$) to $-4.7\%$ (Bin~10, $\bar{z} = 1.85$). 
Since $C_\ell \propto W^2$, the power spectrum deficit ranges
from approximately $-16\%$ at low redshift to $-9\%$ at high
redshift.

This pattern has four distinguishing properties:
(i)~negative at all redshifts, (ii)~monotonically decreasing
with increasing $z$ (the CDM drain is strongest at late
times), (iii)~smooth in $\ell$ (no oscillatory features),
and (iv)~strictly zero in B-modes. No known $\Lambda$CDM
systematic produces this combination. Photometric redshift
errors generate oscillatory residuals in $\ell$; shear
calibration errors are redshift-independent; intrinsic
alignments contaminate B-modes. The EU signature is
geometrically distinct from all three.

\textit{Spectral crossing.}~A further discriminant emerges from the
scale dependence of the $C_\ell$ ratio. At low multipoles
($\ell \lesssim 255$), the geometric kernel bias
dominates: the EU lensing kernels are ${\sim}\,6\%$ weaker
than their $\Lambda$CDM counterparts, suppressing
$C_\ell^{\rm EU}/C_\ell^{\Lambda{\rm CDM}}$ below unity.
At high multipoles ($\ell \gtrsim 255$), the EU matter
power spectrum amplitude takes over: $\sigma_8^{\rm EU} =
0.827$ exceeds $\sigma_8^{\Lambda{\rm CDM}} = 0.811$,
injecting more small-scale power and driving the ratio
above unity. The crossover at $\ell \approx 255$ is a
shape-dependent prediction that does not depend on
the overall normalization of $C_\ell$ --- it depends on
the interplay between a geometric effect (kernel) and
a dynamical effect ($P(k)$ amplitude). Stage-III surveys,
with $\ell_{\rm max} \lesssim 300$, probe only the
suppressed regime. Euclid, reaching $\ell \sim 5000$
across 10 tomographic bins, will map both regimes and
measure the crossover scale directly.

\textit{Pipeline-inferred $\hat{S}_8$.}~The Euclid CUB prediction for the
pipeline-inferred parameter is $\hat{S}_8 = 0.801$. With
Euclid's projected uncertainty $\sigma(S_8) \approx 0.01$
and the Planck uncertainty $\sigma = 0.013$ combined in
quadrature, this lies ${\sim}\,1.9\sigma$ below the Planck
$\Lambda$CDM value ($S_8 = 0.832$). The EU predicts that
this persistent deficit will be interpreted as a growth
tension within $\Lambda$CDM --- but will resolve to
$< 1\sigma$ when the EU lensing kernel is adopted.

\textit{Alcock--Paczy\'nski anomaly.}~The EU expansion
history modifies the radial and transverse distance
scales measured by Euclid's spectroscopic survey
(NISP grism, $0.9 < z < 1.8$). The Alcock--Paczy\'nski
parameters~\cite{alcock_paczynski1979} deviate from
unity:
\begin{equation}
\alpha_\parallel(z{=}1) = 1.012\,,\qquad
\alpha_\perp(z{=}1) = 0.998\,.
\label{eq:ap}
\end{equation}
The radial stretching ($\alpha_\parallel > 1$ by $1.2\%$)
arises because $H_{\rm EU}(z) \neq H_{\Lambda{\rm CDM}}(z)$;
the transverse parameter $\alpha_\perp$ exhibits a sign
change at $z \approx 1.2$, coinciding with the crossover
in the Limber mapping shift
$\Delta k/k = (\chi_{\Lambda{\rm CDM}} - \chi_{\rm EU})/
\chi_{\rm EU}$.
This test is purely spectroscopic and shares no nuisance
parameters with the photometric $3 \times 2$pt analysis.

\textit{$3 \times 2$pt probe decomposition.}~The remaining
legs of the $3 \times 2$pt analysis carry
independent EU signatures. The galaxy clustering kernel
$q(\chi) = b(z)\,n(z)\,H(z)/c$ is biased by
$\Delta H/H$, which ranges from $+0.7\%$ at $z = 0.3$
to $-1.7\%$ at $z = 1.85$ --- changing sign at
$z \approx 0.5$ as the GKI boost
(Sec.~\ref{sec:gki}) gives way to the freezing of the
drain at high redshift. This bias varies
\textit{within} each tomographic bin: $H(z)$ is not
constant across the bin width, so the per-bin galaxy
bias $b_i$ (a free nuisance parameter, constant within
each bin) cannot absorb the signal. The mismatch between
the intra-bin $H(z)$ shape in EU versus $\Lambda$CDM is
a non-degenerate signature.

Galaxy-galaxy lensing $C_\ell^{\delta\gamma}$ inherits
bias from all three physical sources simultaneously:
the lensing kernel deficit ($-5\%$ to $-8\%$), the
clustering kernel shift ($\pm\!2\%$), and the Limber
mapping displacement ($\pm\!1\%$). Their combination
produces a total bias of $-6.0\%$ at $z = 0.3$ to
$-7.1\%$ at $z = 0.96$ --- making galaxy-galaxy lensing
the single most sensitive Euclid observable to the EU
modification. This sensitivity arises because the cross-
correlation probes both the matter distribution (lensing)
and the galaxy field (clustering) through kernels that
are \textit{both} shifted by the altered expansion
history.

The nuisance-free statistic $E_G$~\cite{zhang2007},
defined as the ratio of lensing-galaxy to galaxy-galaxy
angular spectra, cancels galaxy bias $b$ algebraically.
In $\Lambda$CDM, $E_G = \Omega_m/f(z)$, where $f(z)$ is
the linear growth rate. In the EU, both $\Omega_m(z)$ and
$f(z)$ are modified by the drain. The resulting bias
ranges from $-4.7\%$ at $z = 0.3$ (where the CDM deficit
is largest) to $-0.02\%$ at $z = 1.85$ (where the drain
freezes and the EU converges to $\Lambda$CDM). This
convergence is itself a prediction: $E_G(z)$ should
approach the $\Lambda$CDM value at $z > z_{\rm trans}$,
providing a direct measurement of the evaporation onset
redshift.

% ----------------------------------------------------------------
% X.B — LSST Y1
% ----------------------------------------------------------------
\subsection{LSST Y1 weak lensing}
\label{sec:lsst}

The Vera C.~Rubin Observatory Legacy Survey of Space
and Time~\cite{lsst_srd2018} provides a ground-based,
Southern-hemisphere complement to Euclid's Northern
photometric survey. Its Year~1 analysis will employ
five tomographic bins ($z_{\rm edges} = 0.20, 0.43,
0.63, 0.90, 1.30, 3.00$) following the DESC Science
Requirements Document~\cite{mandelbaum2018}, with an
effective source density $n_{\rm eff} = 10\;\mathrm{arcmin}^{-2}$
and intrinsic ellipticity dispersion $\sigma_e = 0.26$.

\textit{The kernel bias gradient.}~Applying the EU lensing
kernel to the LSST $n(z)$ distribution yields a
tomographic gradient in the kernel bias:
\begin{equation}
\Delta W_i/W_i = \{-7.9\%,\; -7.3\%,\; -6.7\%,\; -5.9\%,\; -4.8\%\}\,,
\label{eq:lsst_gradient}
\end{equation}
for bins $i = 1$ through $5$ ($\bar{z} = 0.35$ to $1.75$).
The $40\%$ decrease from Bin~1 to Bin~5 is a direct
imprint of the CDM drain freezing at high redshift:
shallow bins, dominated by galaxies in the active-drain
epoch ($z \lesssim 2$), exhibit the full kernel deficit;
deep bins, weighted toward the frozen regime, are less
affected. This monotonic gradient is unique to the EU
--- no other extension of $\Lambda$CDM produces this
specific $z$-dependent pattern.

The same gradient in Euclid's 10 bins ($-8.1\%$ to
$-4.7\%$) confirms that the signal is physical, not
instrumental. Euclid operates in space at optical-to-NIR
wavelengths in the Northern hemisphere; LSST operates
from the ground at optical wavelengths in the Southern
hemisphere. Independent detection of the same gradient
across both surveys would constitute strong evidence
for a cosmological origin.

\textit{The false internal tension.}~The kernel bias
gradient has a practical consequence for $\Lambda$CDM
analyses. A per-bin $\hat{S}_8$ fit to LSST data in
the EU universe will yield systematically higher
$\hat{S}_8$ values in deep bins (where the kernel bias
is weaker) than in shallow bins (where it is stronger).
This bin-dependent trend will appear as an internal
inconsistency when analyzed under $\Lambda$CDM. The EU
predicts that it will be attributed to photo-$z$
systematic errors at high redshift --- but will resolve
when the CDM drain is included in the analysis kernel.

\textit{Pipeline-inferred $\hat{S}_8$.}~The kernel bias
gradient is shown in Fig.~\ref{fig:kernel_gradient}.
The LSST CUB
prediction is $\hat{S}_8 = 0.801$, consistent with the
Euclid prediction ($\hat{S}_8 = 0.801$). Both forecasts
derive from the same EU background cosmology; the
difference in $n(z)$, sky coverage, and systematics
modeling produces a $< 0.1\%$ shift in the predicted
value. The convergence of two independent Stage-IV
pipelines on the same $\hat{S}_8$ is a consistency
check of the EU prediction.

\begin{figure}[pos=!htbp]
\includegraphics[width=\columnwidth]{fig5_kernel_gradient.pdf}
\caption{Kernel bias gradient $\Delta W_i / W_i$ as a function of
mean tomographic redshift for Euclid DR1 (10 bins, circles) and
LSST Y1 (5 bins, squares). The $40\%$ decrease from
$\bar{z} \approx 0.3$ to $\bar{z} \approx 1.8$ traces the
CDM drain freezing at high redshift.}
\label{fig:kernel_gradient}
\end{figure}

% ----------------------------------------------------------------
% X.C — DESI and the CPL mirage
% ----------------------------------------------------------------
\subsection{Distance probes and the equation-of-state mirage}
\label{sec:cpl_mirage}

The Dark Energy Spectroscopic Instrument~\cite{desi_dr2_2025}
measures baryon acoustic oscillation (BAO) distances $D_H/r_d$
and $D_M/r_d$ across seven tracer populations spanning
$0.1 < z < 2.5$. The EU predicts systematic deviations from
the $\Lambda$CDM baseline that change sign with redshift.

\textit{BAO distance residuals.}~At low redshift ($z \lesssim 0.5$),
the CDM drain increases $H(z)$ through the GKI mechanism
(Sec.~\ref{sec:gki}), yielding $D_H^{\rm EU} < D_H^{\Lambda\rm CDM}$
and $D_M^{\rm EU} < D_M^{\Lambda\rm CDM}$. At $z \gtrsim 0.5$,
the drain freezes and $H_{\rm EU}(z) < H_{\Lambda\rm CDM}(z)$
(because the pre-drain $\Omega_{\rm cdm}$ was larger), reversing
the sign. The $D_H$ crossover occurs at $z \approx 0.5$; the
$D_M$ crossover at $z \approx 1.2$. At $z = 2.33$ (Ly$\alpha$),
the EU predicts $\Delta D_H/D_H = +1.7\%$ and
$\Delta D_M/D_M = +0.6\%$. This sign structure is consistent
with the DESI DR2 finding that low-redshift BAO measurements
fall above the Planck $\Lambda$CDM prediction~\cite{desi_dr2_2025}.

\textit{The CPL mirage.}~The vacuum equation of state is $w_\Lambda = -1$
exactly (Sec.~\ref{sec:stability}), but a CPL fitter
applied to EU-generated distances reveals a
quantitative mirage. We fit a CPL template
$w(a) = w_0 + w_a(1{-}a)$~\cite{chevallier2001, linder2003}
to the EU-generated $D_H(z)$ and $D_M(z)$ at 37 redshift
points via Nelder-Mead minimization. Two configurations
are tested.

\textit{Fit~A} ($\Omega_m$ free):
the fitter recovers $\Omega_m = 0.295$, $w_0 = -1.006$,
$w_a = -0.12$, with $\chi^2 = 2 \times 10^{-5}$. The
pipeline correctly identifies the EU as a cosmological
constant: when allowed to adjust $\Omega_m$ to the
post-drain value ($\Omega_m \approx 0.29$), the CPL
template finds $w_0 \approx -1$ and negligible $w_a$.

\textit{Fit~B} ($\Omega_m$ fixed to CMB prior):
fixing $\Omega_m = 0.3153$ (Planck 2018), the fitter
yields
\begin{equation}
w_0 = -0.855\,,\qquad w_a = -0.81\,.
\label{eq:cpl_mirage}
\end{equation}
The $\chi^2$ degrades to $5.2 \times 10^{-4}$ but remains
small. The CMB prior forces the fitter to accommodate the
drain-reduced distances by pushing $w_0$ above $-1$.

The physical mechanism is transparent: the CDM drain
reduces $\Omega_m$ from $0.315$ (CMB epoch) to $0.288$
(today). When a pipeline anchors $\Omega_m$ to the
CMB value, the excess matter relative to the true
late-time density must be compensated by an apparent
change in the dark energy equation of state. The
CPL template absorbs this mismatch as $w_0 > -1$
(quintessence-like). The apparent ``dynamical dark
energy'' is a prior artifact of the CDM mass
reassignment.

\textit{Comparison with DESI DR2.}~The DESI collaboration
reports~\cite{desi_dr2_2025}
\begin{equation}
w_0 = -0.838 \pm 0.055\,,\quad
w_a = -0.62^{+0.22}_{-0.19}\,,
\label{eq:desi_w0}
\end{equation}
for the combination DESI~BAO + Planck~CMB +
Pantheon+ (Eq.~7 of Ref.~\cite{desi_dr2_2025}),
corresponding to a $2.8\sigma$ preference for evolving
dark energy. The EU Fit~B prediction
($w_0 = -0.855$) lies $0.3\sigma$ from this
central value. The $w_a$ residual ($-0.81$ versus
$-0.62$) reflects the additional constraining power
of the supernova likelihood, which is absent in our
geometric-only fit.

\textit{$\Omega_m$ scan.}~To demonstrate that the
mirage is a monotonic function of the $\Omega_m$
prior, we repeat Fit~B across $\Omega_m \in [0.27, 0.33]$.
The resulting $w_0(\Omega_m)$ curve is strictly increasing:
$w_0 = -1.04$ at $\Omega_m = 0.288$ (the EU post-drain
value) and $w_0 = -0.86$ at $\Omega_m = 0.315$ (Planck).
The crossover $w_0 = -1$ occurs at $\Omega_m \approx 0.297$,
precisely between the EU and $\Lambda$CDM values.
Higher $\Omega_m$ priors yield progressively stronger
quintessence-like signals ($w_0 > -1$), as the fitter
weakens dark energy to compensate for the excess matter.
Conversely, priors below the crossover yield mild
phantom-like values ($w_0 < -1$), as the fitter
strengthens dark energy to compensate for the deficit
of high-redshift matter. The mirage is
driven by $\Delta\Omega_m = \Omega_{m,\rm CMB} -
\Omega_{m,\rm local}$.

This result connects to the ``mirage'' dark energy
class identified in the DESI extended
analysis~\cite{desi_extended_de_2025},
where models along $w_a \approx -3.66(1 + w_0)$
capture the full DE phenomenology with a single
degree of freedom. The authors note that ``the
exact physical mechanism for such rapid emergence
of dark energy remains unclear.'' The CDM drain
provides this mechanism
(Fig.~\ref{fig:cpl_mirage}): it reduces the late-time
matter density, forcing a quintessence-like
compensation when the pipeline assumes a constant
$\Omega_m$.
The same mechanism operates in the growth data:
the $\Lambda$CDM joint constraint
$S_8 = 0.812$ coincides with the EU prediction
$S_8 = 0.811$ because a constant-$\Omega_m$ fit
must approximate the dynamical $\Omega_m(z)$
trajectory of an evaporating universe
(Sec.~\ref{sec:s8_resolution}).

\textit{Growth versus geometry.}~The EU growth
suppression $g = D_{\rm EU}/D_{\Lambda\rm CDM} = 0.993$
is below $1\%$, well within current error bars on
$f\sigma_8$. The DESI Full-Shape analysis confirms
general relativity at the $\mu_0 = 0.11 \pm 0.45$
level~\cite{desi_fullshape_2024}. The EU is consistent
with this null result: the drain modifies
\textit{geometry} (distances, kernels, AP parameters)
but leaves \textit{growth} essentially unchanged.
This asymmetry is the model's distinctive signature
--- it separates the EU from modified-gravity theories,
which generically alter both growth and geometry at
comparable levels.

\begin{figure}[pos=!htbp]
\includegraphics[width=\columnwidth]{fig6_cpl_mirage.pdf}
\caption{The CPL mirage: apparent $w_0$ recovered by a $w$CDM
pipeline as a function of the forced $\Omega_m$ prior. The EU
predicts $w_0 = -0.855$ at $\Omega_m = 0.315$ (red circle),
consistent with the DESI DR2 measurement
$w_0 = -0.838 \pm 0.055$ (blue band) at $0.3\sigma$.
Crossover $w_0 = -1$ at $\Omega_m \approx 0.297$
(green diamond).}
\label{fig:cpl_mirage}
\end{figure}

% ----------------------------------------------------------------
% X.D — Prediction summary and timeline
% ----------------------------------------------------------------
\subsection{Prediction summary and timeline}
\label{sec:timeline}

Table~\ref{tab:predictions} collects ten falsifiable
predictions of the EU model, each fully specified
before encountering the test data.
Prediction~P6, the CPL mirage, is already consistent with
DESI DR2 at $0.3\sigma$.

\begin{table*}[t]
\caption{Falsifiable predictions of the Evaporating Universe.
All quantities are computed from the baseline EU posteriors
and the $N$-body calibrated power spectrum, with zero
additional free parameters. P6 is already tested against
DESI DR2 data.}
\label{tab:predictions}
\begin{tabular}{clllc}
 & Prediction & Observable & Magnitude & Data \\
\hline
P1 & $C_\ell$ suppression (WL)
  & $\Delta C_\ell/C_\ell$
  & $-9\%$ to $-16\%$
  & Euclid DR1 \\
P2 & $\hat{S}_8$ below Planck
  & $\hat{S}_8$
  & $0.801$ ($1.9\sigma$ below $\Lambda$CDM)
  & Euclid DR1 \\
P3 & AP radial stretching
  & $\alpha_\parallel$
  & $1.012$ at $z{=}1$
  & Euclid NISP \\
P4 & $C_\ell$ spectral crossing
  & $C_\ell^{\rm EU}/C_\ell^{\Lambda}$
  & crossover at $\ell \approx 255$
  & Euclid DR1 \\
P5 & Probe inconsistency (GC vs WL)
  & $(w_0, \Omega_m)$ split
  & resolves with CDM drain
  & Euclid DR1 \\
P6 & CPL mirage ($w_0 > -1$)
  & $w_0$
  & $-0.855$ (with CMB prior)
  & \textbf{DESI DR2 (\checkmark)} \\
P7 & LSST $\hat{S}_8$ below Planck
  & $\hat{S}_8$
  & $0.801 \pm 0.012$
  & LSST Y1 \\
P8 & LSST kernel bias gradient
  & $\Delta W/W$ vs $z$
  & $-7.9\%$ to $-4.8\%$
  & LSST Y1 \\
P9 & $w_0$ stabilization
  & $w_0$
  & $\approx -0.85$ (not $\to -1$)
  & DESI DR3 \\
P10 & BAO distance excess at $z > 2$
  & $D_H/r_d,\; D_M/r_d$
  & $+1.7\%,\; +0.6\%$ at $z = 2.33$
  & DESI DR3 \\
\end{tabular}
\end{table*}

The timeline is compressed. Euclid DR1
($\sim\!2500\;\text{deg}^2$, expected 2026) tests
predictions P1--P5 simultaneously. DESI DR2 (2025) has
already provided a $0.3\sigma$-consistent measurement of
P6. LSST Y1 (expected 2027) tests P7 and P8 from the
Southern hemisphere with an independent photometric
pipeline; detection of the same kernel bias gradient
reported by Euclid in the Northern hemisphere would
confirm the cosmological origin of the signal. DESI
DR3 (2027) will test P9 and P10 with improved BAO
statistics and additional tracer populations.

The three programs --- Euclid (space, photometric +
spectroscopic, North), LSST (ground, photometric, South),
and DESI (ground, spectroscopic, all-sky) --- employ
independent instruments, pipelines, and sky coverage.
If all ten predictions are confirmed, the EU
framework will have passed tests spanning photometric
lensing, spectroscopic clustering, BAO distance
measurements, and the dark energy equation of state
--- with zero fitted dark-sector parameters.

% ================================================================
% XI. DISCUSSION
% ================================================================
\section{Discussion}
\label{sec:discussion}

% ----------------------------------------------------------------
% XI.A — Model comparison
% ----------------------------------------------------------------
\subsection{Model comparison}
\label{sec:model_comparison}

Di~Valentino et al.~\cite{divalentino2021} catalog roughly
100 beyond-$\Lambda$CDM proposals aimed at the Hubble or
growth tensions, spanning ten broad categories. We compare
the EU against the strongest representative of each
category in Table~\ref{tab:model_comparison}.

\begin{table}[t]
\caption{Model comparison. $\Delta k$: extra free
parameters. Level: 0\,=\,parametrization,
1\,=\,phenomenological, 2\,=\,form derived / values
fitted, 3\,=\,form and values derived. Tensions
from published posteriors vs~SH0ES and
weak-lensing surveys.
$^\dagger$With SH0ES and KiDS-1000 in fit.}
\label{tab:model_comparison}
\begin{tabular}{p{1.8cm}ccccl}
Model & $\Delta k$ & $H_0$ & $S_8$ & Level & Ref. \\
\hline
$\Lambda$CDM
  & 0 & $5.8\sigma$ & ${\sim}\,2.6\sigma$ & --- &
  \cite{planck2018} \\
EDE
  & 2--3 & ${\sim}\,2\sigma$ &
  ${\sim}\,3\sigma\,\uparrow$ & 2 &
  \cite{poulin2019_ede} \\
$\Lambda_s$CDM
  & 1 & ${\lesssim}\,1\sigma^{\dagger}$ &
  ${\lesssim}\,1\sigma^{\dagger}$ & 1 &
  \cite{akarsu2023} \\
RVM
  & 1 & ${\sim}\,4\sigma$ &
  ${\sim}\,1.5\sigma$ & 2 &
  \cite{sola2017_rvm} \\
IDE ($Q \propto \rho_{\rm de}$)
  & 1--2 & ${\sim}\,2.6\sigma$ & reduces & 1 &
  \cite{divalentino2020} \\
$w_0 w_a$CDM
  & 2 & --- & marginal & 0 &
  \cite{desi_dr2_2025} \\
\textbf{EU}
  & \textbf{0} & $\mathbf{0.52\sigma}$ &
  $\mathbf{0.58\sigma}$ & \textbf{3}
  & This work \\
\end{tabular}
\end{table}

A useful organizing principle is the
\textit{derivation hierarchy} of model parameters:

\begin{itemize}
\item \textbf{Level~0 (parametrization):}
  The functional form is postulated without physical
  motivation. The CPL parametrization
  $w(a) = w_0 + w_a(1-a)$ is the standard example.

\item \textbf{Level~1 (phenomenological):}
  A coupling form is postulated
  (e.g., $Q = \xi H\rhoc$ or $Q = \xi H\rho_{\rm de}$)
  with numerical coefficients fitted to data. This
  describes the IDE models of
  Di~Valentino et al.~\cite{divalentino2020},
  decaying dark matter~\cite{vattis2019},
  and $\Lambda_s$CDM~\cite{akarsu2023}.

\item \textbf{Level~2 (form derived, values fitted):}
  The coupling form is motivated by quantum field theory,
  but numerical parameters remain free.
  EDE~\cite{poulin2019_ede} derives an axion-like potential
  from the string landscape; the Running Vacuum
  Model~\cite{sola2017_rvm} derives
  $\Lambda(H) = c_0 + \nu H^2$ from QFT in curved
  spacetime. In both cases, at least one coefficient
  ($f_{\rm EDE}$, $z_c$, or $\nu$) is fitted by MCMC.

\item \textbf{Level~3 (form and values derived):}
  Both the functional form and all numerical parameters
  are determined by the theoretical framework, with
  $\Delta k = 0$. The EU is the only published model
  in this category.
\end{itemize}

The closest competitor in derivation rigor is the
Running Vacuum Model of Sol\`a
Peracaula~\cite{sola2017_rvm}. Both the RVM and the EU
invoke QFT to modify the dark sector; both produce
a coupling of the form $Q \propto H\rhoc$. The
difference is surgical: the RVM running coefficient
$\nu$ is \textit{motivated} by renormalization group
arguments but \textit{fitted} by MCMC, placing it at
Level~2. The EU coupling
$\lambda\varepsilon_{\rm IR} = (2/3)({\pi}/{4})/\sqrt{108\pi}$
is \textit{calculated} from the Gribov--Zwanziger
effective action and \textit{confirmed} by MCMC
(Sec.~\ref{sec:concordance}), placing it at Level~3.

Phenomenological IDE models ($Q = \xi H\rhoc$)
and decaying dark matter~\cite{vattis2019}
are the Level~1 counterparts of the EU:
they postulate the same coupling form that the EU
derives. The EU thus \textit{provides the microscopic
mechanism} that the IDE literature assumes.

The DESI DR2 measurement of $w_0 > -1$ at
${\sim}\,2.8\sigma$~\cite{desi_dr2_2025} is often
interpreted as evidence for dynamical dark energy.
Within the EU, the vacuum equation of state is
$w_\Lambda = -1$ exactly (Sec.~\ref{sec:stability}),
but an observer fitting CPL templates to
EU-generated distances recovers $w_0 = -1.006$
when $\Omega_m$ is free (the EU \textit{is} a
cosmological constant) and $w_0 = -0.855$ when
$\Omega_m$ is fixed to the Planck CMB prior
(Sec.~\ref{sec:cpl_mirage}). The latter lies
$0.3\sigma$ from the DESI central value
($w_0 = -0.838 \pm 0.055$), confirming that
the DESI signal is quantitatively consistent with
a CDM drain artifact.


% ----------------------------------------------------------------
% XI.B — Theoretical structure and limitations
% ----------------------------------------------------------------
\subsection{Theoretical structure and limitations}
\label{sec:theory_limits}

The EU rests on a physical distinction between the 4D
spacetime bulk and the 3D spatial hypersurface $\Sigma$.
While the ADM decomposition is conventionally treated as
a gauge choice, the Gribov--Zwanziger ghost
condensate~\cite{gribov1978,zwanziger1989,vandersickel2012}
is a non-perturbative phenomenon defined on the 4D
manifold that affects $\Sigma$ through a calculable
barrier. The 10 independent components of $g_{\mu\nu}$
and the six independent components of $\gamma_{ij}$ define
two physically distinct domains; the four constrained
degrees of freedom (lapse and shift, removed by the
Hamiltonian and momentum constraints) constitute the
interface between them. The requirement that all vacuum
fluctuations respect spherical isotropy --- the same
geometric principle that fixes the factor $4\pi$ in
Gauss's law --- recurs at every stage of the EU
derivation: in the heat-kernel prefactor
$(4\pi)^{-d/2}|_{d=4}$ that sets $z_{\rm trans}$,
in the loop measure $(4\pi)^{-1}$ that enters
$\varepsilon_{\rm bare}$, and in the CKN volume
$(4\pi/3)R_H^3$ that establishes the UV cutoff.
This geometric coherence is the reason
$\Delta k = 0$: no free parameters survive when
isotropy fully determines the structure at every scale.

The triple convergence on $\lambda = 2/3$
(Sec.~\ref{sec:lambda}) and the Slavnov--Taylor
protection of this value to all orders
(Eq.~\ref{eq:ward}) ensure that the effective coupling
$\varepsilon_{\rm IR} = (\pi/4)\,\varepsilon_{\rm bare}$
contains zero adjustable coefficients.

\textit{Why baryons are immune.}~The democratic coupling
derived in Sec.~\ref{sec:eps} requires the interacting
field to be a gauge singlet under the Ashtekar SU(2)
connection: only singlet states satisfy the
Wigner--Eckart condition for uniform coupling to all
27 ghost channels. Baryonic matter carries Standard
Model gauge charges
($\mathrm{SU}(3)_{\rm QCD} \times \mathrm{U}(1)_{\rm EM}$)
and transforms non-trivially under the full gauge group.
Its internal quantum numbers break the democratic
coupling condition, excluding it from the condensate
interaction. This selection rule is structural: it is
the negative corollary of the same representation
theory that determines the CDM coupling amplitude
$\varepsilon_{\rm bare}$ in Eq.~(\ref{eq:epsbare}).

\textit{Limitations.}~Several aspects of the present
analysis require qualification.
(i)~The growth tension is resolved to
$\sim\!1\sigma$ using the CUB forward model
(Sec.~\ref{sec:kernel_bias}), but the baryonic
feedback correction relies on the ratio
$P_{\rm EU}/P_{\Lambda{\rm CDM}}$ from a
gravity-only $N$-body simulation
(Sec.~\ref{sec:nbody}). A hydrodynamic simulation
incorporating the CDM drain alongside baryonic
processes would sharpen this estimate.
(ii)~The local kinematic imprint (LKI) accounts for
$\sim\!70\%$ of the $H_0$ resolution.
The KBC void parameters ($\delta \approx -0.30$,
$R \sim 300$~Mpc) are observational inputs, not
EU-specific predictions; the EU permits and
quantifies the void boost, but does not create
the void.
(iii)~The framework has been validated with a single
Boltzmann code implementation (CLASS-EU).
Independent reproduction --- using, e.g., CAMB with
equivalent patches --- is a natural next step
following publication. The full MCMC infrastructure
(CLASS patches, Cobaya configuration files,
Gadget-4 scripts) will be made publicly available
to facilitate this.
(iv)~The CDM drain modifies $H(z)$ at the percent
level, by design. Differential tests therefore
require high-precision data; the predictions of
Sec.~\ref{sec:predictions} are calibrated to the
sensitivity of Euclid, LSST, and DESI DR3.

% ----------------------------------------------------------------
% XI.C — Outlook
% ----------------------------------------------------------------
\subsection{Outlook}
\label{sec:outlook}

A concrete prediction emerges from comparing the
$N$-body and CUB forward models across angular scales.
At $\ell \lesssim 300$ --- the regime probed by
Stage-III surveys --- the gravity-only simulation
yields $\hat{S}_{8,\rm Nbody} < \hat{S}_{8,\rm CUB}$,
because CDM evaporation makes halos genuinely lighter
(``anemic'' suppression). At
$\ell \lesssim 5000$ --- accessible to Euclid ---
the inequality reverses:
$\hat{S}_{8,\rm Nbody} > \hat{S}_{8,\rm CUB}$,
because the DM-only simulation lacks baryonic blowout
and retains artificially dense cores.
The crossover occurs at $\ell \approx 300$, precisely
where current surveys place their scale cuts.
The amplitude of the reversal
($\Delta S_8 \approx 0.09$) quantifies the baryonic
feedback correction and constitutes a falsifiable
prediction for Euclid's extended multipole coverage.

DESI DR3 (2027) will provide a sharper test of the
equation-of-state duality (Sec.~\ref{sec:cpl_mirage}).
A probe-split analysis --- fitting $w_0$ and $w_a$
separately from galaxy clustering and weak lensing ---
should reveal inconsistent $(w_0, \Omega_m)$ values
under $\Lambda$CDM, resolving to concordance when the
CDM drain is included.

The cosmographic bias identified in
Sec.~\ref{sec:pipeline_bias} --- contamination of
$M_B$ by the KBC void outflow --- will receive a
complete forensic decomposition in Paper~II. That
analysis will combine Cosmicflows-4 peculiar
velocity fields~\cite{tully2023}, a joint
CMB$\,{+}\,$BAO$\,{+}\,$SNe fit with $M_B$ free, and
the statistical amplification of deep voids in the EU
power spectrum ($\sigma(R)$ enhanced by $3.5\%$
relative to $\Lambda$CDM, rendering voids of
$\delta > 0.3$ an order of magnitude more probable).
The objective is to complete the parallel established
in Sec.~\ref{sec:s8_resolution}: the growth tension
is resolved by the lensing kernel bias; the Hubble
tension, by the cosmographic bias. Both are
consequences of analyzing an evaporating universe
with a non-evaporating template.

% ================================================================
% XII. CONCLUSIONS
% ================================================================
\section{Conclusions}
\label{sec:conclusions}

The Evaporating Universe is an interacting dark energy model
in which a Gribov--Zwanziger ghost condensate mediates a
continuous transfer of energy from CDM to the vacuum. All
five model parameters are derived from EFT in curved
spacetime; none are fitted to data ($\Delta k = 0$).
The principal findings are as follows.

\begin{enumerate}

\item \textit{Framework.}~The coupling
$Q = -\lambda\,\varepsilon(z)\,H\,\rhoc$ is the unique
leading-order interaction consistent with FLRW symmetry
and Fermi's golden rule. The dissipation fraction
$\lambda = 2/3$, logistic exponent $b = 19/36$, transition
redshift $z_{\rm trans} \approx 5.986$, bare amplitude
$\varepsilon_{\rm bare} \approx 0.0543$, and effective
coupling $\varepsilon_{\rm IR} \approx 0.0426$ are
determined by the spatial dimensionality $n = 3$ and the
topological structure of the Gribov horizon. The
Cancellation Theorem guarantees $w_\Lambda = -1$ exactly,
eliminating the doom factor instability that plagues
generic IDE models.

\item \textit{Hubble tension.}~The CDM drain raises the
global expansion rate to
$H_0^{\rm GKI} = 68.89 \pm 0.28$~km\,s$^{-1}$\,Mpc$^{-1}$.
The KBC void outflow adds a local kinematic imprint,
yielding
$H_0^{\rm LKI} = 72.71 \pm 0.30$~km\,s$^{-1}$\,Mpc$^{-1}$
--- within $0.52\sigma$ of
SH0ES~\cite{breuval2024}. Adding the SH0ES likelihood
to the MCMC shifts the posterior by $0.25\sigma$,
confirming that the local measurement is redundant
rather than constraining.

\item \textit{Growth tension.}~Weak-lensing surveys
infer $S_8$ using a $\Lambda$CDM lensing kernel that
systematically underestimates the amplitude in an
evaporating universe. Forward-modeling the EU kernel
yields $\hat{S}_8 = 0.794$ for
DES-Y3~\cite{des2022}, reducing the tension from
${\sim}\,2.6\sigma$ to $1.03\sigma$ (kernel-corrected).
Including DES-Y3 data directly in the MCMC yields
$S_8 = 0.815 \pm 0.007$ --- within $0.58\sigma$ of the
blind EU prediction ($0.811$).

\item \textit{Validation.}~All six MCMC runs converge at
$R-1 < 0.01$, including the 38-dimensional
EU+$H_0$+$S_8$ configuration. The UV-derived
parameters match the Phase~2 EU posteriors (CMB+BAO+SNe)
to $0.04\sigma$~(Sec.~\ref{sec:concordance}).
A Gadget-4 $N$-body simulation ($1024^3$ particles)
confirms the drain signature in $P(k)$.
Model selection yields
$\Delta\text{AIC} = -6.88$ in the fair-fight
configuration ($\Delta k = 0$, both models
void-corrected).

\item \textit{Falsifiability.}~Ten zero-parameter
predictions target Euclid DR1, LSST Y1, and DESI DR3.
Prediction P6 --- the CPL mirage
($w_0 = -0.855$ when $\Omega_m$ is fixed to the CMB
prior) --- already matches the DESI DR2 measurement
($w_0 = -0.838 \pm 0.055$) at $0.3\sigma$
(Sec.~\ref{sec:cpl_mirage}). The remaining nine span
$C_\ell$ suppression ($-9\%$ to $-16\%$),
$\hat{S}_8 = 0.801$ ($1.9\sigma$ below Planck under
$\Lambda$CDM),
Alcock--Paczy\'nski stretching
($\alpha_\parallel = 1.012$ at $z = 1$),
LSST kernel bias gradient
($-7.9\%$ to $-4.8\%$ across five bins), and
BAO distance excess at $z > 2$.
If any prediction fails, the framework is
falsified --- because all derive from the same
modified $H(z)$.

\end{enumerate}

The Evaporating Universe resolves both cosmological
tensions simultaneously, with zero additional free
parameters --- not by fitting the data, but by
predicting it.


% ================================================================
% ACKNOWLEDGMENTS
% ================================================================
\section*{Acknowledgments}

This work would not have been possible without the open-source tools and
public data sets maintained by the cosmology community. I gratefully
acknowledge the developers of
\textsc{class}~\cite{blas2011},
\textsc{Cobaya}~\cite{torrado2021},
\textsc{CosmoLike}~\cite{krause2021},
\textsc{Gadget-4}~\cite{gadget4},
\textsc{monofonIC}~\cite{monofonic},
\textsc{Pylians}~\cite{pylians},
and \textsc{HMCode}~\cite{mead2021}
for making state-of-the-art numerical cosmology accessible to
independent researchers. I am equally indebted to the
Planck~\cite{planck2018}, DES~\cite{des2022},
KiDS~\cite{heymans2021},
HSC~\cite{li2023_hsc},
SH0ES~\cite{breuval2024}, Pantheon$+$~\cite{scolnic2022},
and DESI~\cite{desi_dr2_2025}
collaborations for their commitment to public data releases.

All code, data, and supplementary material for this work are
publicly available~\cite{eu_github}.

I am deeply grateful to my mother, Nilceia Camarinho, for her
unwavering dedication to my education and moral foundation,
filling every role to ensure my growth. A special thanks to my
wife, Rayane Brum, my perfect life partner, whose unconditional
support and nurturing presence allowed me to fully explore my
potential while keeping me tethered to reality---preventing me
from sharing the solitary fate of Jos\'e Arcadio
Buend\'ia. To them, I owe the path that led to this work.

% ================================================================
% DATA AVAILABILITY
% ================================================================
\section*{DATA AVAILABILITY}

The modified \textsc{class} source code, analysis notebooks,
MCMC chains, and $N$-body simulation outputs that support the
findings of this article are openly available~\cite{eu_github}.

% ================================================================
% DECLARATION OF COMPETING INTERESTS
% ================================================================
\section*{Declaration of competing interests}

The author declares that he has no known competing financial
interests or personal relationships that could have appeared
to influence the work reported in this article.

% ================================================================
% FUNDING
% ================================================================
\section*{Funding}

This research did not receive any specific grant from funding
agencies in the public, commercial, or not-for-profit sectors.

% ================================================================
% APPENDICES
% ================================================================
\appendix

% ================================================================
% APPENDIX A --- DETAILED DERIVATIONS
% ================================================================
\section{Detailed derivations supporting Sec.~\ref{sec:params}}
\label{app:derivations}

% ----------------------------------------------------------------
% A.1 --- TT Projection
% ----------------------------------------------------------------
\subsection{Ghost-graviton vertex and TT projection}
\label{app:tt_vertex}

The ghost-graviton vertex in $d = 4$ linearized quantum gravity
reads~\cite{codello2009}
\begin{align}
V^\mu_{\;\nu;\rho\sigma}(p,q) = \tfrac{1}{2}\big[
  &\delta^\mu_\rho\, q_\sigma\, p_\nu
  + \delta^\mu_\sigma\, q_\rho\, p_\nu
  + \delta_{\nu\rho}\, p_\sigma\, q^\mu
  \notag \\
  &+ \delta_{\nu\sigma}\, p_\rho\, q^\mu
  - \delta_{\rho\sigma}(q^\mu p_\nu + p^\mu q_\nu)
  \notag \\
  &- \delta^\mu_\nu(p_\rho q_\sigma + p_\sigma q_\rho
  - \delta_{\rho\sigma}\,p{\cdot}q)\big],
\label{eq:vertex}
\end{align}
where $p$ is the ghost momentum and $q = p + k$ the anti-ghost
momentum.

The York TT projector on symmetric rank-2
tensors~\cite{york1973} is
\begin{equation}
\Pi^{\rm TT}_{\rho\sigma,\lambda\tau}(k) =
  \tfrac{1}{2}(P_{\rho\lambda}P_{\sigma\tau}
  + P_{\rho\tau}P_{\sigma\lambda})
  - \tfrac{1}{d{-}1}P_{\rho\sigma}P_{\lambda\tau},
\label{eq:projector}
\end{equation}
where $P_{\mu\nu} = \delta_{\mu\nu} - k_\mu k_\nu/k^2$ is the
transverse projector, satisfying idempotency
($\Pi^{\rm TT}\cdot\Pi^{\rm TT} = \Pi^{\rm TT}$), transversality
($k^\rho\Pi^{\rm TT}_{\rho\sigma,\lambda\tau} = 0$), and
tracelessness
($\delta^{\rho\sigma}\Pi^{\rm TT}_{\rho\sigma,\lambda\tau} = 0$).

The TT trace ratio
$b = {\rm Tr}(V{\cdot}\Pi^{\rm TT}{\cdot}V)/
     {\rm Tr}(V{\cdot}\mathbb{I}{\cdot}V)$
is evaluated using irreducible exact rational arithmetic, ensuring
zero floating-point contamination. Four independent momentum
configurations yield:

\begin{table}[pos=!htbp]
\caption{Verification of $b = 19/36$ across four momentum
  configurations. All computations use exact rational arithmetic.}
\label{tab:b_verification}
\begin{tabular}{lccc}
Configuration & ${\rm Tr}_{\rm TT}/p^2$ &
  ${\rm Tr}_{\rm full}/p^2$ & Ratio \\
\hline
$p \perp k$, axes (1,2) & $19/3$ & $12$ & $19/36$ \\
$p \perp k$, axes (2,3) & $19/3$ & $12$ & $19/36$ \\
$p \perp k$, axes (4,1) & $19/3$ & $12$ & $19/36$ \\
$|p|{=}|k|{=}\sqrt{2}$  & $38/3$ & $24$ & $19/36$ \\
\end{tabular}
\end{table}

\noindent
The ratio is momentum-independent, as required for a well-defined
projector fraction. The identity on symmetric tensors entering the
denominator is
$\mathbb{I}_{\rho\sigma,\lambda\tau} =
  \frac{1}{2}(\delta_{\rho\lambda}\delta_{\sigma\tau}
  + \delta_{\rho\tau}\delta_{\sigma\lambda})$.

% ----------------------------------------------------------------
% A.2 --- z_trans Sensitivity
% ----------------------------------------------------------------
\subsection{Condensation threshold sensitivity}
\label{app:ztrans_sensitivity}

The transition redshift depends weakly on $h = H_0/100$ through
$\Om = (\omega_b + \omega_{\rm cdm})/h^2$, where $\omega_b = 0.02237$
and $\omega_{\rm cdm} = 0.1200$ are direct CMB
observables~\cite{planck2018}. For $h \in [0.660, 0.700]$:

\begin{table}[pos=!htbp]
\caption{Sensitivity of $\ztrans$ to the reduced Hubble constant $h$.
  The total variation is $\Delta\ztrans \approx 0.40$ ($\sim 6.7\%$).}
\label{tab:ztrans_sensitivity}
\begin{tabular}{cccl}
$h$ & $\Om$ & $\ztrans$ & Source \\
\hline
$0.660$ & $0.328$ & $5.847$ & Lower bound \\
$0.6736$ & $0.315$ & $5.986$ & Planck ($\Lambda$CDM) \\
$0.690$ & $0.300$ & $6.150$ & EU upper estimate \\
$0.700$ & $0.292$ & $6.248$ & Upper bound \\
\end{tabular}
\end{table}

\noindent
Higher-order heat-kernel corrections are of order
$(H/\MPl)^2 \sim 10^{-122}$, confirming that the leading-order
threshold is exact for all practical purposes.

% ----------------------------------------------------------------
% A.3 --- Gribov Integral
% ----------------------------------------------------------------
\subsection{Gribov propagator integral}
\label{app:gribov_integral}

The screening ratio $I_{\rm IR}/I_{\rm UV}$ is evaluated on the
two-dimensional transverse momentum plane. This requires
justification beyond the channel counting of
Sec.~\ref{sec:eps}: eliminating the longitudinal
polarization ($\lambda = 2/3$) reduces the number of active
\textit{channels} but does not by itself reduce the dimensionality of
the momentum integral --- a photon has two polarizations yet
integrates over $d^3k$.

The 2D measure follows from Zwanziger's horizon dominance
theorem~\cite{zwanziger1991}: in the thermodynamic limit
$V \to \infty$, the functional measure concentrates on the Gribov
horizon $\partial\Omega$ --- a codimension-1 surface in configuration
space --- for the same reason that the volume of a high-dimensional
sphere concentrates at its surface. The CDM$\,\to\,$vacuum transition
crosses this boundary (Sec.~\ref{sec:gz}), so by Fermi's
golden rule the density of accessible final states is defined on
$\partial\Omega$, yielding the transverse measure $k\,dk/(2\pi)$.

An independent route converges on the same result: the CKN
holographic bound (Sec.~\ref{sec:ztrans}, Link~2) establishes that
the effective degrees of freedom in a causal volume are governed by
its bounding surface, not its bulk. Both the Zwanziger argument
(field-space boundary) and the CKN argument (geometric boundary)
select the 2D measure.

\textit{UV reference} (free massive propagator,
$D_{\rm UV} = 1/(k^2 + \gamma^2)$):
\begin{equation}
I_{\rm UV} = \frac{1}{2\pi}\int_0^\infty
  \frac{k\,dk}{(k^2 + \gamma^2)^2}
  = \frac{1}{4\pi\gamma^2}\,.
\label{eq:IUV}
\end{equation}

\textit{IR integral} (Gribov propagator,
$D_{\rm GZ} = k^2/(k^4 + \gamma^4)$):
\begin{equation}
I_{\rm IR} = \frac{1}{2\pi}\int_0^\infty
  \frac{k^5\,dk}{(k^4 + \gamma^4)^2}\,.
\label{eq:IIR}
\end{equation}
Substituting $u = k^2$, then $u = \gamma^2\tan\theta$:
\begin{align}
I_{\rm IR}
  &= \frac{1}{4\pi}\int_0^\infty
    \frac{u^2\,du}{(u^2 + \gamma^4)^2}
  \notag \\
  &= \frac{1}{4\pi\gamma^2}\int_0^{\pi/2}
    \sin^2\!\theta\;d\theta
  = \frac{1}{4\pi\gamma^2}\cdot\frac{\pi}{4}
  = \frac{1}{16\gamma^2}\,.
\label{eq:IIR_result}
\end{align}
The ratio is
\begin{equation}
\frac{I_{\rm IR}}{I_{\rm UV}}
  = \frac{1/(16\gamma^2)}{1/(4\pi\gamma^2)}
  = \frac{\pi}{4}\,,
\label{eq:ratio}
\end{equation}
as can be verified by direct evaluation. The Gribov mass $\gamma$
cancels in the ratio, yielding a universal screening factor.


% ================================================================
% ================================================================


% APPENDIX B --- SUPPORTING ANALYSES FOR SEC. V
% ================================================================
\section{Supporting analyses for Sec.~\ref{sec:mcmc}}
\label{app:free_params}

\subsection{Perturbation null test}
\label{app:null_test}

\begin{table*}[pos=!htbp]
\caption{Perturbation null test: Phase~1 (perturbation
source terms off) versus Phase~2 (source terms on).
The EU-blind configuration frees
$\varepsilon_{\rm IR}$, $b$, $z_{\rm trans}$; the
theory-fixed configuration holds them at the
framework-derived values. All posterior shifts are
$< 0.03\sigma$.}
\label{tab:null_test}
\begin{tabular}{lcccc}
 & \multicolumn{2}{c}{EU-blind ($\Delta k{=}3$)}
 & \multicolumn{2}{c}{EU ($\Delta k{=}0$)} \\
 & Phase~1 (off) & Phase~2 (on)
 & Phase~1 (off) & Phase~2 (on) \\
\hline
\multicolumn{5}{c}{\textit{Sampled parameters}} \\
$\omega_b$
  & $0.02231 \pm 0.00012$ & $0.02231 \pm 0.00012$
  & $0.02218 \pm 0.00012$ & $0.02218 \pm 0.00013$ \\
$\omega_{\rm cdm}$
  & $0.11777 \pm 0.00061$ & $0.11778 \pm 0.00062$
  & $0.11927 \pm 0.00064$ & $0.11926 \pm 0.00063$ \\
$100\theta_s$
  & $1.04191 \pm 0.00023$ & $1.04191 \pm 0.00023$
  & $1.04180 \pm 0.00023$ & $1.04180 \pm 0.00023$ \\
$\tau_{\rm reio}$
  & $0.059 \pm 0.007$ & $0.059 \pm 0.007$
  & $0.056 \pm 0.007$ & $0.056 \pm 0.007$ \\
$\ln(10^{10}A_s)$
  & $3.046 \pm 0.014$ & $3.047 \pm 0.014$
  & $3.044 \pm 0.014$ & $3.044 \pm 0.014$ \\
$n_s$
  & $0.9676 \pm 0.0034$ & $0.9676 \pm 0.0034$
  & $0.9637 \pm 0.0035$ & $0.9637 \pm 0.0034$ \\
\hline
\multicolumn{5}{c}{\textit{Derived parameters}} \\
$H_0^{\rm GKI}$ (km\,s$^{-1}$\,Mpc$^{-1}$)
  & $68.13 \pm 0.28$ & $68.13 \pm 0.28$
  & $68.88 \pm 0.28$ & $68.89 \pm 0.28$ \\
$\sigma_8$
  & $0.805 \pm 0.006$ & $0.806 \pm 0.006$
  & $0.828 \pm 0.006$ & $0.827 \pm 0.006$ \\
$S_8$
  & $0.810 \pm 0.008$ & $0.810 \pm 0.008$
  & $0.811 \pm 0.008$ & $0.811 \pm 0.008$ \\
$\Omega_m$
  & $0.303 \pm 0.004$ & $0.303 \pm 0.004$
  & $0.288 \pm 0.003$ & $0.288 \pm 0.003$ \\
\hline
\multicolumn{5}{c}{\textit{Diagnostics}} \\
$\chi^2_{\rm total}$
  & $12{,}394.7$ & $12{,}394.4$
  & $12{,}396.2$ & $12{,}396.3$ \\
$(R-1)_{\rm max}$
  & $0.003$ & $0.003$
  & $0.004$ & $0.005$ \\
\end{tabular}
\end{table*}

Table~\ref{tab:null_test} presents the full posterior
comparison between Phase~1 (perturbation source terms off)
and Phase~2 (source terms on) for both the EU-blind
($\Delta k = 3$) and theory-fixed ($\Delta k = 0$)
configurations. All parameter shifts are below
$0.03\sigma$, confirming the cancellation theorem of
Sec.~\ref{sec:stability}: the perturbation source terms
have no measurable effect on the posterior.

\subsection{Free-parameter runs}
\label{app:free_runs}

The EU-blind configuration samples three additional
parameters with uniform priors:
$\varepsilon_{\rm IR} \in [0.001, 0.200]$ (theory: $0.04264$),
$b \in [0.10, 2.00]$ (theory: $19/36$), and
$z_{\rm trans} \in [2.0, 12.0]$ (theory: $5.986$).
Figure~\ref{fig:corner_comparison} shows the resulting
posterior comparison. The three-parameter degeneracy
and its statistical consequences are discussed in
Sec.~\ref{sec:posteriors}; the theory-fixed configuration
is preferred by $\Delta\text{AIC} = -4.1$.


\begin{figure*}[pos=!htbp]
\includegraphics[width=0.85\textwidth]{fig7_corner_comparison.pdf}
\caption{Posterior comparison between the theory-fixed
configuration (blue) and the EU-blind free-parameter
configuration (green). Diagonal panels show
1D marginalized posteriors; off-diagonal panels show
68\% and 95\% credible regions. Red dotted lines mark
the framework-derived values.}
\label{fig:corner_comparison}
\end{figure*}

\clearpage

% ================================================================
% DECLARATION OF GENERATIVE AI
% ================================================================
\section*{Declaration of generative AI and AI-assisted technologies in the manuscript preparation process}

During the preparation of this work, the author used an AI
coding assistant in order to support manuscript editing,
structural organization, editorial compliance verification,
and computational code development (CLASS-EU solver patches,
Cobaya MCMC configuration files, and Gadget-4 $N$-body
simulation scripts). After using this tool, the author
reviewed and edited all content as needed and takes full
responsibility for the content of the published article.

% To print the credit authorship contribution details
\printcredits

% BIBLIOGRAPHY
% ================================================================
\bibliographystyle{cas-model2-names}
\bibliography{eu_paper1}

\end{document}
